% ============================================================
%  PROGRAMMER'S GUIDE
%  Bug-Bite Classification via Tropical Cyclone Transfer Learning
%  Braude College of Engineering — Software Engineering Capstone
%
%  Usage in Overleaf:
%    \appendix
%    % ============================================================
%  PROGRAMMER'S GUIDE
%  Bug-Bite Classification via Tropical Cyclone Transfer Learning
%  Braude College of Engineering — Software Engineering Capstone
%
%  Usage in Overleaf:
%    \appendix
%    % ============================================================
%  PROGRAMMER'S GUIDE
%  Bug-Bite Classification via Tropical Cyclone Transfer Learning
%  Braude College of Engineering — Software Engineering Capstone
%
%  Usage in Overleaf:
%    \appendix
%    % ============================================================
%  PROGRAMMER'S GUIDE
%  Bug-Bite Classification via Tropical Cyclone Transfer Learning
%  Braude College of Engineering — Software Engineering Capstone
%
%  Usage in Overleaf:
%    \appendix
%    \input{programmer_guide}
%
%  Or compile standalone:
%    Uncomment the preamble block below, and add \end{document} at the end.
% ============================================================

% --- Standalone preamble (remove if including in main document) ----------
\documentclass[12pt,a4paper]{article}
\usepackage[utf8]{inputenc}
\usepackage[T1]{fontenc}
\usepackage{geometry}
\geometry{margin=2.5cm}
\usepackage{hyperref}
\usepackage{listings}
\usepackage{xcolor}
\usepackage{booktabs}
\usepackage{array}
\usepackage{enumitem}
\usepackage{parskip}
\lstset{basicstyle=\ttfamily\small, breaklines=true, frame=single,
        backgroundcolor=\color{gray!8}}
\begin{document}
\title{Programmer's Guide\\
       \large Bug-Bite Classification via Tropical Cyclone Transfer Learning\\
       Braude College of Engineering --- Software Engineering Capstone}
\author{Ishay Yulzary, Nathan Trostianitser}
\date{}
\maketitle
\tableofcontents
\newpage
% -------------------------------------------------------------------------

\section{Programmer's Guide}
\label{app:programmer-guide}

This appendix is addressed to a software engineer who wishes to understand the
internals of the system, reproduce the experimental results, extend the research,
or advance the prototype toward a deployable product.
All source code lives in the repository root unless otherwise noted.

% ============================================================
\subsection{Research Overview}
\label{sec:pg-overview}
% ============================================================

\subsubsection{Problem and Motivation}

Classifying bug-bite skin images by causative species is clinically useful but
technically difficult: inter-class visual similarity is high, intra-class variation
is large, and labelled medical image datasets are small. Standard fine-tuning from
ImageNet leaves the backbone biased toward object-centric features that are of
limited relevance to dermatological macro-photography.

\subsubsection{Core Hypothesis}

This project investigates whether inserting an intermediate pre-training step on
tropical cyclone (TC) infrared satellite imagery improves backbone representations
for the bug-bite task. The conjecture is that cyclone satellite images and skin
macro-photographs share structural visual properties: radial textures, soft
gradients, and fine-grained pattern differences between severity categories that
are invisible to coarse object detectors. A backbone that has already learned to
discriminate these properties at the cyclone level may require less adaptation to
do so at the skin level.

\subsubsection{Transfer Classification Pipeline}

Two training pipelines are compared:

\begin{description}
    \item[Control] ImageNet pre-training $\rightarrow$ bug-bite fine-tuning
          (direct, single-stage transfer).
    \item[TC] ImageNet pre-training $\rightarrow$ cyclone intensity pre-training
          $\rightarrow$ bug-bite fine-tuning (two-stage transfer).
\end{description}

The cyclone pre-training stage is a full classification task: the model learns to
distinguish five cyclone intensity categories derived from the Saffir-Simpson
scale. The resulting backbone weights are then transplanted into a fresh
bug-bite classification head, and the whole network is fine-tuned.

\subsubsection{Experimental Design}

Three backbone architectures are evaluated under both pipelines:

\begin{center}
\begin{tabular}{lll}
\toprule
\textbf{Architecture} & \textbf{timm identifier} & \textbf{Input size} \\
\midrule
ConvNeXt-Tiny  & \texttt{convnext\_tiny.fb\_in22k\_ft\_in1k} & 256$\times$256 \\
DenseNet-121   & \texttt{densenet121.ra\_in1k}                & 256$\times$256 \\
InceptionV3    & \texttt{inception\_v3.tv\_in1k}              & 299$\times$299 \\
\bottomrule
\end{tabular}
\end{center}

The cyclone pre-training hyperparameters are swept across the following grid,
yielding 63 configurations in total (3 label-smoothing values $\times$ 3 learning
rates $\times$ 7 random seeds):

\begin{center}
\begin{tabular}{ll}
\toprule
\textbf{Parameter} & \textbf{Values} \\
\midrule
Cyclone label smoothing (\texttt{ls}) & 0.0, 0.1, 0.2 \\
Cyclone phase-2 learning rate (\texttt{lr}) & $10^{-6}$, $5\times10^{-6}$, $10^{-5}$ \\
Random seeds & 0, 1, 2, 3, 4, 5, 6 \\
\bottomrule
\end{tabular}
\end{center}

Bug-bite fine-tuning hyperparameters are fixed across all configurations:
label smoothing 0.0, phase-2 learning rate $5\times10^{-6}$.

% ============================================================
\subsection{Repository Structure}
\label{sec:pg-repo}
% ============================================================

\begin{description}
    \item[\texttt{notebooks/Multiclass\_Classification.ipynb}]
          The primary experimental notebook. Covers all four pipeline stages:
          control training, TC pre-training, feature map and GradCAM evaluation,
          and the full XAI suite. Intended for interactive exploration or
          one-shot reproduction on a GPU instance.

    \item[\texttt{Stacked\_Model\_Colab.ipynb}]
          Self-contained inference and visualisation notebook designed for
          Google Colab free tier. Loads trained weights from Google Drive,
          runs ensemble inference on a single image, and produces side-by-side
          Control vs.\ TC GradCAM and LIME outputs.

    \item[\texttt{scripts/}]
          Standalone Python scripts, one per pipeline stage. Intended for
          execution on a headless server via \texttt{run\_experiments.py}.
          \begin{itemize}
              \item \texttt{shared.py} --- shared imports, dataset paths,
                    metrics helpers, and GradCAM utilities used by all four scripts.
              \item \texttt{01\_train\_control.py} --- trains the Control
                    (ImageNet $\rightarrow$ bug-bite) models.
              \item \texttt{02\_train\_tc.py} --- runs cyclone pre-training then
                    TC $\rightarrow$ bug-bite fine-tuning.
              \item \texttt{03\_evaluate.py} --- generates feature map comparisons
                    and GradCAM overlays.
              \item \texttt{04\_xai.py} --- computes ensemble metrics, silhouette
                    scores, t-SNE projections, LIME, SHAP, and DiCE outputs.
          \end{itemize}

    \item[\texttt{miscellaneous\_code/pytorch\_utils.py}]
          Custom PyTorch training library imported by both the notebook and the
          scripts. See Section~\ref{sec:pg-pytorch-utils} for full documentation.

    \item[\texttt{miscellaneous\_code/run\_experiments.py}]
          Automated hyperparameter sweep runner. See
          Section~\ref{sec:pg-sweep} for full documentation.

    \item[\texttt{miscellaneous\_code/fetch\_cyclone\_dataset.py}]
          Downloads and labels the HURSAT-B1 cyclone dataset from NOAA.
          See Section~\ref{sec:pg-data-cyclone}.

    \item[\texttt{miscellaneous\_code/cyclone\_preprocessing.py}]
          Removes grid-line artefacts introduced by the HURSAT-B1 satellite
          image tiling process before training.

    \item[\texttt{Bug-Data/Multiclass\_Bug\_data/}]
          Bug-bite dataset, structured as \texttt{train/} and \texttt{val/}
          subdirectories, each containing one subdirectory per class.

    \item[\texttt{results/}]
          Auto-generated by the sweep runner. Contains checkpoints, per-run
          metric JSONs, XAI outputs, and plots.

    \item[\texttt{deploy/}]
          Legacy Streamlit demo (binary classifier, outdated). Not part of
          the active pipeline; retained for reference only.
\end{description}

% ============================================================
\subsection{Core Training Library: \texttt{pytorch\_utils.py}}
\label{sec:pg-pytorch-utils}
% ============================================================

This file contains all reusable PyTorch infrastructure. It is imported by
\texttt{run\_experiments.py}, by the four pipeline scripts in \texttt{scripts/},
and by the main notebook. Engineers extending the system should understand each
public function before modifying the training loop.

\subsubsection{\texttt{get\_pytorch\_loaders}}

\begin{lstlisting}[language=Python]
get_pytorch_loaders(train_dir, val_dir,
                    img_size=310, batch_size=4, augment=False)
\end{lstlisting}

Returns \texttt{(train\_loader, val\_loader, class\_names)} using
\texttt{torchvision.datasets.ImageFolder}. The \texttt{augment} parameter
accepts three values:

\begin{description}
    \item[\texttt{False}] No augmentation. Used for validation and for
          bug-bite fine-tuning in the sweep (the bug-bite dataset is small;
          augmentation during fine-tuning degraded performance in preliminary
          experiments).
    \item[\texttt{True}] Light augmentation: horizontal flip, $\pm10^{\circ}$
          rotation. Used for moderate augmentation scenarios.
    \item[\texttt{'strong'}] Aggressive augmentation designed for cyclone
          satellite imagery: random resized crop (scale 0.35--1.0), full
          $360^{\circ}$ rotation, vertical flip, colour jitter, random greyscale,
          Gaussian blur (kernel 5, $\sigma \in [0.5, 3.0]$), and random erasing.
          The full rotation and vertical flip are valid for cyclone images because
          satellite imagery is radially symmetric with respect to the storm centre.
\end{description}

All loaders use ImageNet normalisation
(\texttt{mean=[0.485, 0.456, 0.406]}, \texttt{std=[0.229, 0.224, 0.225]}).

\subsubsection{\texttt{train\_pytorch\_model}}

\begin{lstlisting}[language=Python]
train_pytorch_model(model, train_loader, val_loader, device,
                    phase1_epochs=5, phase2_epochs=20, patience=5,
                    save_path=None,
                    phase1_batch_size=None, phase2_batch_size=None,
                    label_smoothing=0.0, auto_class_weights=False,
                    phase1_lr=1e-4, phase2_lr=5e-6)
\end{lstlisting}

Implements the standard two-phase fine-tuning strategy used throughout the
project:

\begin{description}
    \item[Phase 1 (head-only)] All backbone parameters are frozen.
          Only the classification head (detected via \texttt{model.get\_classifier()})
          is trained for \texttt{phase1\_epochs} epochs at \texttt{phase1\_lr}
          with AdamW ($\lambda = 0.01$).
    \item[Phase 2 (full fine-tune)] All parameters are unfrozen and the
          entire network is trained at \texttt{phase2\_lr}. Gradient checkpointing
          is enabled where the model supports it.
\end{description}

Both phases use:
\begin{itemize}
    \item Early stopping on validation loss with patience \texttt{patience}.
    \item Automatic Mixed Precision (AMP) with \texttt{bfloat16} when a CUDA
          device is available.
    \item \texttt{CrossEntropyLoss} with optional label smoothing.
    \item Patience counters that reset independently between phases,
          because phase 2 typically causes an initial loss spike after
          unfreezing the backbone.
\end{itemize}

The best checkpoint (lowest validation loss) is saved to \texttt{save\_path}
and automatically reloaded into the model before the function returns.

\subsubsection{\texttt{get\_backbone\_state\_dict}}

\begin{lstlisting}[language=Python]
get_backbone_state_dict(state_dict)
\end{lstlisting}

Strips all classifier head keys from a state dict, enabling backbone weight
transfer between tasks with different numbers of output classes. Keys whose
names start with \texttt{head.}, \texttt{fc.}, or \texttt{classifier.} are
removed. This function is the mechanism by which the cyclone-pre-trained
backbone is transplanted into a fresh bug-bite classification head.

\subsubsection{Supporting Functions}

\begin{description}
    \item[\texttt{get\_pt\_probs\_preds(model, loader, device)}]
          Returns \texttt{(true\_labels, argmax\_preds, softmax\_probs)} as
          Python lists. Used for evaluation after training.
    \item[\texttt{evaluate\_pytorch\_model(model, val\_loader, class\_names, device)}]
          Prints a full \texttt{sklearn} classification report and renders a
          confusion matrix.
    \item[\texttt{plot\_pytorch\_history(history)}]
          Renders train/validation accuracy and loss curves from the history
          dict returned by \texttt{train\_pytorch\_model}.
\end{description}

% ============================================================
\subsection{Dataset Acquisition}
\label{sec:pg-data}
% ============================================================

\subsubsection{Bug-Bite Dataset}
\label{sec:pg-data-bug}

The bug-bite dataset contains images of five species categories:
\texttt{ants}, \texttt{bed\_bugs}, \texttt{mosquitos}, \texttt{spiders},
and \texttt{ticks\_fleas}.
It is organised as \texttt{train/\{class\}/} and \texttt{val/\{class\}/}
subdirectories inside \texttt{Bug-Data/Multiclass\_Bug\_data/}.

The dataset is available via two channels:
\begin{enumerate}
    \item \textbf{Repository:} the \texttt{Bug-Data/} directory is included in
          the repository and can be used directly.
    \item \textbf{Pre-processed archive:} a curated and pre-processed version,
          including augmented samples used in the published experiments, is
          available on request from\\\texttt{ishay.yulzary@gmail.com}.
\end{enumerate}

Dataset paths are configured at the top of \texttt{miscellaneous\_code/run\_experiments.py}
(\texttt{BUG\_TRAIN}, \texttt{BUG\_VAL}) and in \texttt{scripts/shared.py}.
They default to WSL native paths (\texttt{/home/test/bug\_data/\{train,val\}}),
which must be updated for other environments.

\subsubsection{Cyclone Dataset}
\label{sec:pg-data-cyclone}

The cyclone dataset is derived from two public sources:
\begin{itemize}
    \item \textbf{HURSAT-B1} (Hurricane Satellite, Band 1): infrared brightness
          temperature satellite imagery at $8\,\text{km}$ resolution,
          hosted by NOAA.
    \item \textbf{IBTrACS} (International Best Track Archive for Climate
          Stewardship): per-observation wind speed records used to assign
          intensity labels to each satellite frame.
\end{itemize}

Five intensity categories are defined, mapped to the Saffir-Simpson scale:

\begin{center}
\begin{tabular}{lll}
\toprule
\textbf{Category label} & \textbf{Wind speed (kt)} & \textbf{Saffir-Simpson equivalent} \\
\midrule
\texttt{cat1\_tropical\_depression} & 0--33   & Tropical Depression \\
\texttt{cat2\_tropical\_storm}      & 34--63  & Tropical Storm \\
\texttt{cat3\_hurricane\_cat12}     & 64--95  & Hurricane Cat 1--2 \\
\texttt{cat4\_hurricane\_cat34}     & 96--136 & Hurricane Cat 3--4 \\
\texttt{cat5\_hurricane\_cat5}      & 137+    & Hurricane Cat 5 \\
\bottomrule
\end{tabular}
\end{center}

Images cover the years 1995--2016. Each frame is clipped to the brightness
temperature range $[180, 300]\,\text{K}$, inverted (so cold cloud tops appear
bright), and saved as a $224\times224$ PNG.

\subsubsection*{Fetching the dataset}

\begin{lstlisting}[language=bash]
# Full dataset: 4000 images per category (20000 total, ~4h on a stable connection)
python miscellaneous_code/fetch_cyclone_dataset.py

# Quick trial run: 5 images per category
python miscellaneous_code/fetch_cyclone_dataset.py --trial

# Custom target size and worker count
python miscellaneous_code/fetch_cyclone_dataset.py --target 2000 --workers 8
\end{lstlisting}

The script is resumable: interrupted runs pick up from where they stopped
using \\\texttt{cyclone\_dataset/download\_log.json}. The output directory
defaults to \texttt{cyclone\_dataset/} in the working directory.

After fetching, the dataset must be split into train/val and placed at the
paths expected by the sweep runner. The pre-built split (500 val images per
category) is available from \texttt{ishay.yulzary@gmail.com}.

% ============================================================
\subsection{Training Pipeline: Individual Scripts}
\label{sec:pg-scripts}
% ============================================================

The four scripts in \texttt{scripts/} mirror the four sections of the main
notebook. They accept the same arguments and produce the same outputs, but
are designed for unattended server execution.

\subsubsection{\texttt{01\_train\_control.py}}

Trains the three Control models (ConvNeXt-Tiny, DenseNet-121, InceptionV3)
for one seed. Saves checkpoint files at the paths supplied via CLI arguments.

\begin{lstlisting}[language=bash]
python scripts/01_train_control.py \
    --run-id s0_ls0.10_lr5.0e-06 \
    --seed 0 \
    --ctrl-convnext  results/checkpoints/control_seed_0/control_convnext.pt \
    --ctrl-densenet  results/checkpoints/control_seed_0/control_densenet.pt \
    --ctrl-inception results/checkpoints/control_seed_0/control_inception.pt
\end{lstlisting}

\subsubsection{\texttt{02\_train\_tc.py}}

For one configuration, runs cyclone pre-training on all three architectures
and then fine-tunes each cyclone backbone on the bug-bite task.
Requires both cyclone checkpoint paths (output) and TC checkpoint paths (output).

\begin{lstlisting}[language=bash]
python scripts/02_train_tc.py \
    --run-id s0_ls0.10_lr5.0e-06 \
    --seed 0 \
    --label-smoothing 0.1 \
    --phase2-lr 5e-6 \
    --cyc-convnext  results/checkpoints/config_000/cyclone_convnext.pt \
    --cyc-densenet  results/checkpoints/config_000/cyclone_densenet.pt \
    --cyc-inception results/checkpoints/config_000/cyclone_inception.pt \
    --tc-convnext   results/checkpoints/config_000/tc_bug_convnext.pt \
    --tc-densenet   results/checkpoints/config_000/tc_bug_densenet.pt \
    --tc-inception  results/checkpoints/config_000/tc_bug_inception.pt
\end{lstlisting}

\subsubsection{\texttt{03\_evaluate.py}}

Generates feature map visualisations and GradCAM overlays comparing the
Control backbone, the TC cyclone backbone (before bug-bite fine-tuning),
and the TC fine-tuned model. Outputs are saved to
\texttt{results/feature\_maps/\{run-id\}/} and
\texttt{results/plots/\{run-id\}/}.

\subsubsection{\texttt{04\_xai.py}}

Computes the full XAI suite for one configuration:
ensemble accuracy (Control vs.\ TC), silhouette scores on backbone embeddings,
t-SNE projections, LIME superpixel explanations, SHAP gradient attributions,
and DiCE counterfactual explanations. Outputs are saved to
\texttt{results/xai/\{run-id\}/}.
Optional XAI libraries (\texttt{shap}, \texttt{lime}, \texttt{dice-ml}) are
checked at runtime; cells that depend on missing libraries are skipped
gracefully.

% ============================================================
\subsection{Automated Sweep Runner: \texttt{run\_experiments.py}}
\label{sec:pg-sweep}
% ============================================================

\texttt{miscellaneous\_code/run\_experiments.py} is the primary entry point for
training. It iterates over the full 63-configuration hyperparameter grid, calls
the four pipeline scripts in sequence for each configuration, and manages
checkpointing, state, and results atomically.

\subsubsection{Basic Usage}

\begin{lstlisting}[language=bash]
# Run all pending configurations
python miscellaneous_code/run_experiments.py

# Check progress without running
python miscellaneous_code/run_experiments.py --status

# Wipe state and results (checkpoints are preserved)
python miscellaneous_code/run_experiments.py --reset
\end{lstlisting}

\subsubsection{Configuration Grid}

The grid is defined by three constants at the top of the file:

\begin{lstlisting}[language=Python]
LABEL_SMOOTHING_VALUES = [0.0, 0.1, 0.2]
PHASE2_LR_VALUES       = [1e-6, 5e-6, 1e-5]
N_SEEDS                = 7
\end{lstlisting}

Adding a value to \texttt{LABEL\_SMOOTHING\_VALUES} or \texttt{PHASE2\_LR\_VALUES}
and re-running automatically appends the new configurations to the existing
state file without invalidating completed work.

\subsubsection{Pause and Resume}

The sweep can be interrupted and resumed at any point:

\begin{description}
    \item[Graceful pause between models:] create the file
          \texttt{results/PAUSE} (\texttt{touch results/PAUSE}) or press
          Ctrl-C. The runner finishes the current model, saves its checkpoint,
          and exits cleanly.
    \item[Crash recovery:] any configuration marked \texttt{running} in the
          state file (indicating a previous interrupted run) is automatically
          reset to \texttt{pending}. Already-saved checkpoints are reused,
          so only the incomplete model needs to be retrained.
\end{description}

\subsubsection{Output Files}

\begin{description}
    \item[\texttt{results/experiment\_state.json}]
          Tracks the status of every configuration (\texttt{pending},
          \texttt{running}, or \texttt{completed}).
    \item[\texttt{results/experiment\_results.json}]
          Per-model, per-configuration evaluation results: validation accuracy,
          macro F1, precision, recall, and full classification report.
    \item[\texttt{results/checkpoints/config\_\{id\}/}]
          TC cyclone checkpoint (\texttt{cyclone\_\{model\}.pt}) and
          TC bug-bite checkpoint (\texttt{tc\_bug\_\{model\}.pt}) for each
          configuration.
    \item[\texttt{results/checkpoints/control\_seed\_\{s\}/}]
          Control checkpoints, keyed by seed only (the Control pipeline is
          independent of the cyclone hyperparameters).
    \item[\texttt{results/xai/\{run-id\}/}]
          Per-configuration XAI outputs: \texttt{ensemble\_metrics.json},
          \\\texttt{silhouette.json}, t-SNE plots, GradCAM overlays,
          LIME visualisations, SHAP attribution maps.
\end{description}

\subsubsection{Hardware Requirements}

The full 63-configuration sweep was run on an NVIDIA RTX A5000 (24\,GB VRAM).
A GPU with at least 16\,GB VRAM is recommended for training. Estimated runtime
per configuration is 1.5--2 hours on equivalent hardware.
Inference and visualisation (scripts \texttt{03} and \texttt{04}) can be run
on CPU or a smaller GPU.

% ============================================================
\subsection{Experimental Results}
\label{sec:pg-results}
% ============================================================

\subsubsection{Ensemble Accuracy: Control vs.\ TC}

The table below reports the mean ensemble validation accuracy across all seven
seeds for each hyperparameter configuration. The ensemble averages softmax
probabilities across all three architectures before taking the argmax.

\begin{center}
\begin{tabular}{lcrrrr}
\toprule
\textbf{Config (ls, lr)} & \textbf{N seeds} &
\textbf{Control mean} & \textbf{TC mean} & \textbf{$\Delta$} \\
\midrule
ls=0.10, lr=$5\times10^{-6}$ & 7 & 81.83\% & 83.27\% & +1.45 pp \\
ls=0.00, lr=$1\times10^{-6}$ & 7 & 81.83\% & 83.18\% & +1.36 pp \\
ls=0.20, lr=$1\times10^{-6}$ & 7 & 81.83\% & 83.18\% & +1.36 pp \\
ls=0.20, lr=$5\times10^{-6}$ & 7 & 81.83\% & 83.00\% & +1.18 pp \\
ls=0.20, lr=$1\times10^{-5}$ & 7 & 81.83\% & 82.82\% & +0.99 pp \\
ls=0.00, lr=$1\times10^{-5}$ & 7 & 81.83\% & 82.73\% & +0.90 pp \\
ls=0.00, lr=$5\times10^{-6}$ & 7 & 81.83\% & 82.73\% & +0.90 pp \\
ls=0.10, lr=$1\times10^{-6}$ & 7 & 81.83\% & 82.28\% & +0.45 pp \\
ls=0.10, lr=$1\times10^{-5}$ & 7 & 81.83\% & 81.83\% & +0.00 pp \\
\bottomrule
\end{tabular}
\end{center}

TC pre-training matches or exceeds the Control baseline in all nine
configurations. The best single-run result is \textbf{86.08\%} TC ensemble
accuracy (seed 4, ls=0.00, lr=$10^{-5}$), compared to 81.01\% Control for
the same seed and configuration.

\subsubsection{Backbone Embedding Quality: Silhouette Scores}

Silhouette scores are computed on backbone embeddings extracted from the
validation set (before the classification head). Higher scores indicate
better class separation in the feature space.

\begin{center}
\begin{tabular}{lrr}
\toprule
\textbf{Architecture} & \textbf{Control} & \textbf{TC} \\
\midrule
ConvNeXt-Tiny & 0.057 & 0.081 \\
DenseNet-121  & 0.027 & 0.030 \\
InceptionV3   & \text{-}0.003 & \text{-}0.017 \\
\bottomrule
\end{tabular}
\end{center}

(Values shown are from the representative seed 0, ls=0.10, lr=$5\times10^{-6}$
configuration. Per-configuration scores are available in
\texttt{results/xai/\{run-id\}/silhouette.json}.)

% ============================================================
\subsection{XAI Methods}
\label{sec:pg-xai}
% ============================================================

All XAI outputs are generated by \texttt{scripts/04\_xai.py} and by Section 4
of the main notebook.

\begin{description}
    \item[GradCAM] Class Activation Mapping via gradients. Implemented in
          \texttt{scripts/shared.py}. Target layers are architecture-specific:
          the final stage for ConvNeXt, \texttt{norm5} for DenseNet,
          and \texttt{Mixed\_7c} for InceptionV3. Produces side-by-side
          Control vs.\ TC attention overlays for direct visual comparison.

    \item[LIME] Superpixel-based local explanation using
          \texttt{lime.lime\_image.LimeImageExplainer}. The ensemble prediction
          function is used as the black-box model. Top-$N$ segments
          (default 12) are coloured by their contribution sign and weight,
          using a red--blue diverging palette. Computationally expensive at
          high \texttt{num\_samples} values; the notebook default is 1000
          samples per image.

    \item[SHAP] Pixel-level attribution via \texttt{shap.GradientExplainer}.
          A background distribution of 50 validation images is used.
          Attribution maps are averaged across the three ensemble models
          and overlaid on the input image.

    \item[DiCE] Counterfactual explanations in embedding space. Backbone
          embeddings are reduced with PCA (50 components), and a surrogate
          logistic regression classifier is fitted on the validation set.
          DiCE generates counterfactual embedding vectors showing the minimal
          directional change required to flip the predicted class.

    \item[t-SNE] 2D projection of backbone embeddings for all five bug-bite
          classes, displayed as a $3 \times 2$ grid (one row per architecture,
          columns for Control and TC). Used to visualise whether TC pre-training
          produces more separable cluster structure.

    \item[Ensemble comparison] Soft-vote ensemble across all three architectures,
          evaluated on the validation set. Control and TC ensembles are compared
          directly; results saved to \texttt{ensemble\_metrics.json}.
\end{description}

% ============================================================
\subsection{Demo Notebook: \texttt{Stacked\_Model\_Colab.ipynb}}
\label{sec:pg-demo}
% ============================================================

This notebook provides a self-contained inference and visualisation interface
designed to run on Google Colab free tier. It does not require a local GPU or
local dataset.

\paragraph{Prerequisites.}
\begin{enumerate}
    \item A Google account with access to the project Google Drive folder.
          The folder must contain \texttt{project\_book\_results.tar},
          which packages all trained weight files.
    \item The notebook set \texttt{PROJECT\_ROOT} to the correct Drive path
          \\(default: \texttt{/content/drive/MyDrive/Capstone}).
\end{enumerate}

\paragraph{Workflow.}
\begin{enumerate}
    \item Mount Google Drive (Cell 2).
    \item Extract weights from the archive if not already extracted (Cell 3).
    \item Load all six models: three Control and three TC (Cell 5).
    \item Set \texttt{IMG\_PATH} to any JPEG or PNG image (Cell 7).
    \item Run inference: Cell 7 displays prediction probabilities for both
          pipelines side-by-side.
    \item Run GradCAM: Cell 8 renders a $3 \times 2$ overlay grid
          (architectures $\times$ pipelines).
    \item Run LIME: Cell 9 renders superpixel attributions for both pipelines.
\end{enumerate}

Weight file locations within the extracted archive:
\begin{itemize}
    \item Control: \texttt{project\_book\_results/s3\_ls0.20\_lr1.0e-05/}
    \item TC: \texttt{project\_book\_results/s6\_ls0.10\_lr5.0e-06/}
\end{itemize}

% ============================================================
\subsection{Extending the System}
\label{sec:pg-extend}
% ============================================================

\subsubsection{Adding a New Architecture}

\begin{enumerate}
    \item Add an entry to the \texttt{MODELS} list in
          \texttt{miscellaneous\_code/run\_experiments.py}:
          \texttt{(key, timm\_id, img\_size, cyc\_p1\_bs, bug\_bs,
          cyc\_p2\_epochs, bug\_p2\_epochs)}.
    \item Identify the correct GradCAM target layer for the new architecture
          and add a case to the \texttt{\_get\_target} function in
          \texttt{Stacked\_Model\_Colab.ipynb} and in \texttt{scripts/shared.py}.
    \item Re-run the sweep. New configurations are appended automatically;
          existing results are not invalidated.
\end{enumerate}

\subsubsection{Adding a New Bug-Bite Class}

\begin{enumerate}
    \item Add labelled images to \texttt{Bug-Data/Multiclass\_Bug\_data/train/}
          and \texttt{val/} under a new subdirectory named after the class.
    \item Update \texttt{N\_CLASSES} in \texttt{run\_experiments.py} if
          necessary (currently 5; \texttt{timm} models are instantiated
          with this value).
    \item Retrain from scratch. The existing Control and TC checkpoints are
          incompatible with a different number of output classes.
\end{enumerate}

\subsubsection{Adapting to a New Intermediate Domain}

The TC pipeline is not specific to cyclone imagery. Any image classification
task can serve as the intermediate pre-training domain by:
\begin{enumerate}
    \item Placing the new domain dataset at \texttt{TC\_TRAIN} / \texttt{TC\_VAL}
          in \texttt{run\_experiments.py}.
    \item Adjusting the augmentation strategy in \texttt{get\_pytorch\_loaders}
          if the new domain does not share cyclone imagery's symmetry properties.
    \item Updating the \texttt{augment='strong'} branch of
          \texttt{pytorch\_utils.py} accordingly.
\end{enumerate}

\subsubsection{Path Toward Productisation}

The current system is a research prototype. The following steps are required
before deployment:

\begin{enumerate}
    \item \textbf{Model serving:} wrap the ensemble inference logic from
          \texttt{Stacked\_Model\_Colab.ipynb} in a REST API endpoint
          (e.g.\ FastAPI). The core inference function is approximately
          30 lines of PyTorch code.
    \item \textbf{Weight hosting:} move trained weights from Google Drive to
          a persistent object store (e.g.\ AWS S3, GCP Cloud Storage).
    \item \textbf{Input validation:} add image size checks, colour space
          normalisation, and rejection of non-photographic inputs.
    \item \textbf{Dataset expansion:} the current bug-bite dataset is small
          (fewer than 500 images per class). Classification performance will
          improve substantially with a larger, clinically curated dataset.
    \item \textbf{Clinical validation:} model outputs should be reviewed and
          validated by a medical professional before any medical application.
\end{enumerate}

% --- End of appendix (remove if standalone) ---
\end{document}

%
%  Or compile standalone:
%    Uncomment the preamble block below, and add \end{document} at the end.
% ============================================================

% --- Standalone preamble (remove if including in main document) ----------
\documentclass[12pt,a4paper]{article}
\usepackage[utf8]{inputenc}
\usepackage[T1]{fontenc}
\usepackage{geometry}
\geometry{margin=2.5cm}
\usepackage{hyperref}
\usepackage{listings}
\usepackage{xcolor}
\usepackage{booktabs}
\usepackage{array}
\usepackage{enumitem}
\usepackage{parskip}
\lstset{basicstyle=\ttfamily\small, breaklines=true, frame=single,
        backgroundcolor=\color{gray!8}}
\begin{document}
\title{Programmer's Guide\\
       \large Bug-Bite Classification via Tropical Cyclone Transfer Learning\\
       Braude College of Engineering --- Software Engineering Capstone}
\author{Ishay Yulzary, Nathan Trostianitser}
\date{}
\maketitle
\tableofcontents
\newpage
% -------------------------------------------------------------------------

\section{Programmer's Guide}
\label{app:programmer-guide}

This appendix is addressed to a software engineer who wishes to understand the
internals of the system, reproduce the experimental results, extend the research,
or advance the prototype toward a deployable product.
All source code lives in the repository root unless otherwise noted.

% ============================================================
\subsection{Research Overview}
\label{sec:pg-overview}
% ============================================================

\subsubsection{Problem and Motivation}

Classifying bug-bite skin images by causative species is clinically useful but
technically difficult: inter-class visual similarity is high, intra-class variation
is large, and labelled medical image datasets are small. Standard fine-tuning from
ImageNet leaves the backbone biased toward object-centric features that are of
limited relevance to dermatological macro-photography.

\subsubsection{Core Hypothesis}

This project investigates whether inserting an intermediate pre-training step on
tropical cyclone (TC) infrared satellite imagery improves backbone representations
for the bug-bite task. The conjecture is that cyclone satellite images and skin
macro-photographs share structural visual properties: radial textures, soft
gradients, and fine-grained pattern differences between severity categories that
are invisible to coarse object detectors. A backbone that has already learned to
discriminate these properties at the cyclone level may require less adaptation to
do so at the skin level.

\subsubsection{Transfer Classification Pipeline}

Two training pipelines are compared:

\begin{description}
    \item[Control] ImageNet pre-training $\rightarrow$ bug-bite fine-tuning
          (direct, single-stage transfer).
    \item[TC] ImageNet pre-training $\rightarrow$ cyclone intensity pre-training
          $\rightarrow$ bug-bite fine-tuning (two-stage transfer).
\end{description}

The cyclone pre-training stage is a full classification task: the model learns to
distinguish five cyclone intensity categories derived from the Saffir-Simpson
scale. The resulting backbone weights are then transplanted into a fresh
bug-bite classification head, and the whole network is fine-tuned.

\subsubsection{Experimental Design}

Three backbone architectures are evaluated under both pipelines:

\begin{center}
\begin{tabular}{lll}
\toprule
\textbf{Architecture} & \textbf{timm identifier} & \textbf{Input size} \\
\midrule
ConvNeXt-Tiny  & \texttt{convnext\_tiny.fb\_in22k\_ft\_in1k} & 256$\times$256 \\
DenseNet-121   & \texttt{densenet121.ra\_in1k}                & 256$\times$256 \\
InceptionV3    & \texttt{inception\_v3.tv\_in1k}              & 299$\times$299 \\
\bottomrule
\end{tabular}
\end{center}

The cyclone pre-training hyperparameters are swept across the following grid,
yielding 63 configurations in total (3 label-smoothing values $\times$ 3 learning
rates $\times$ 7 random seeds):

\begin{center}
\begin{tabular}{ll}
\toprule
\textbf{Parameter} & \textbf{Values} \\
\midrule
Cyclone label smoothing (\texttt{ls}) & 0.0, 0.1, 0.2 \\
Cyclone phase-2 learning rate (\texttt{lr}) & $10^{-6}$, $5\times10^{-6}$, $10^{-5}$ \\
Random seeds & 0, 1, 2, 3, 4, 5, 6 \\
\bottomrule
\end{tabular}
\end{center}

Bug-bite fine-tuning hyperparameters are fixed across all configurations:
label smoothing 0.0, phase-2 learning rate $5\times10^{-6}$.

% ============================================================
\subsection{Repository Structure}
\label{sec:pg-repo}
% ============================================================

\begin{description}
    \item[\texttt{notebooks/Multiclass\_Classification.ipynb}]
          The primary experimental notebook. Covers all four pipeline stages:
          control training, TC pre-training, feature map and GradCAM evaluation,
          and the full XAI suite. Intended for interactive exploration or
          one-shot reproduction on a GPU instance.

    \item[\texttt{Stacked\_Model\_Colab.ipynb}]
          Self-contained inference and visualisation notebook designed for
          Google Colab free tier. Loads trained weights from Google Drive,
          runs ensemble inference on a single image, and produces side-by-side
          Control vs.\ TC GradCAM and LIME outputs.

    \item[\texttt{scripts/}]
          Standalone Python scripts, one per pipeline stage. Intended for
          execution on a headless server via \texttt{run\_experiments.py}.
          \begin{itemize}
              \item \texttt{shared.py} --- shared imports, dataset paths,
                    metrics helpers, and GradCAM utilities used by all four scripts.
              \item \texttt{01\_train\_control.py} --- trains the Control
                    (ImageNet $\rightarrow$ bug-bite) models.
              \item \texttt{02\_train\_tc.py} --- runs cyclone pre-training then
                    TC $\rightarrow$ bug-bite fine-tuning.
              \item \texttt{03\_evaluate.py} --- generates feature map comparisons
                    and GradCAM overlays.
              \item \texttt{04\_xai.py} --- computes ensemble metrics, silhouette
                    scores, t-SNE projections, LIME, SHAP, and DiCE outputs.
          \end{itemize}

    \item[\texttt{miscellaneous\_code/pytorch\_utils.py}]
          Custom PyTorch training library imported by both the notebook and the
          scripts. See Section~\ref{sec:pg-pytorch-utils} for full documentation.

    \item[\texttt{miscellaneous\_code/run\_experiments.py}]
          Automated hyperparameter sweep runner. See
          Section~\ref{sec:pg-sweep} for full documentation.

    \item[\texttt{miscellaneous\_code/fetch\_cyclone\_dataset.py}]
          Downloads and labels the HURSAT-B1 cyclone dataset from NOAA.
          See Section~\ref{sec:pg-data-cyclone}.

    \item[\texttt{miscellaneous\_code/cyclone\_preprocessing.py}]
          Removes grid-line artefacts introduced by the HURSAT-B1 satellite
          image tiling process before training.

    \item[\texttt{Bug-Data/Multiclass\_Bug\_data/}]
          Bug-bite dataset, structured as \texttt{train/} and \texttt{val/}
          subdirectories, each containing one subdirectory per class.

    \item[\texttt{results/}]
          Auto-generated by the sweep runner. Contains checkpoints, per-run
          metric JSONs, XAI outputs, and plots.

    \item[\texttt{deploy/}]
          Legacy Streamlit demo (binary classifier, outdated). Not part of
          the active pipeline; retained for reference only.
\end{description}

% ============================================================
\subsection{Core Training Library: \texttt{pytorch\_utils.py}}
\label{sec:pg-pytorch-utils}
% ============================================================

This file contains all reusable PyTorch infrastructure. It is imported by
\texttt{run\_experiments.py}, by the four pipeline scripts in \texttt{scripts/},
and by the main notebook. Engineers extending the system should understand each
public function before modifying the training loop.

\subsubsection{\texttt{get\_pytorch\_loaders}}

\begin{lstlisting}[language=Python]
get_pytorch_loaders(train_dir, val_dir,
                    img_size=310, batch_size=4, augment=False)
\end{lstlisting}

Returns \texttt{(train\_loader, val\_loader, class\_names)} using
\texttt{torchvision.datasets.ImageFolder}. The \texttt{augment} parameter
accepts three values:

\begin{description}
    \item[\texttt{False}] No augmentation. Used for validation and for
          bug-bite fine-tuning in the sweep (the bug-bite dataset is small;
          augmentation during fine-tuning degraded performance in preliminary
          experiments).
    \item[\texttt{True}] Light augmentation: horizontal flip, $\pm10^{\circ}$
          rotation. Used for moderate augmentation scenarios.
    \item[\texttt{'strong'}] Aggressive augmentation designed for cyclone
          satellite imagery: random resized crop (scale 0.35--1.0), full
          $360^{\circ}$ rotation, vertical flip, colour jitter, random greyscale,
          Gaussian blur (kernel 5, $\sigma \in [0.5, 3.0]$), and random erasing.
          The full rotation and vertical flip are valid for cyclone images because
          satellite imagery is radially symmetric with respect to the storm centre.
\end{description}

All loaders use ImageNet normalisation
(\texttt{mean=[0.485, 0.456, 0.406]}, \texttt{std=[0.229, 0.224, 0.225]}).

\subsubsection{\texttt{train\_pytorch\_model}}

\begin{lstlisting}[language=Python]
train_pytorch_model(model, train_loader, val_loader, device,
                    phase1_epochs=5, phase2_epochs=20, patience=5,
                    save_path=None,
                    phase1_batch_size=None, phase2_batch_size=None,
                    label_smoothing=0.0, auto_class_weights=False,
                    phase1_lr=1e-4, phase2_lr=5e-6)
\end{lstlisting}

Implements the standard two-phase fine-tuning strategy used throughout the
project:

\begin{description}
    \item[Phase 1 (head-only)] All backbone parameters are frozen.
          Only the classification head (detected via \texttt{model.get\_classifier()})
          is trained for \texttt{phase1\_epochs} epochs at \texttt{phase1\_lr}
          with AdamW ($\lambda = 0.01$).
    \item[Phase 2 (full fine-tune)] All parameters are unfrozen and the
          entire network is trained at \texttt{phase2\_lr}. Gradient checkpointing
          is enabled where the model supports it.
\end{description}

Both phases use:
\begin{itemize}
    \item Early stopping on validation loss with patience \texttt{patience}.
    \item Automatic Mixed Precision (AMP) with \texttt{bfloat16} when a CUDA
          device is available.
    \item \texttt{CrossEntropyLoss} with optional label smoothing.
    \item Patience counters that reset independently between phases,
          because phase 2 typically causes an initial loss spike after
          unfreezing the backbone.
\end{itemize}

The best checkpoint (lowest validation loss) is saved to \texttt{save\_path}
and automatically reloaded into the model before the function returns.

\subsubsection{\texttt{get\_backbone\_state\_dict}}

\begin{lstlisting}[language=Python]
get_backbone_state_dict(state_dict)
\end{lstlisting}

Strips all classifier head keys from a state dict, enabling backbone weight
transfer between tasks with different numbers of output classes. Keys whose
names start with \texttt{head.}, \texttt{fc.}, or \texttt{classifier.} are
removed. This function is the mechanism by which the cyclone-pre-trained
backbone is transplanted into a fresh bug-bite classification head.

\subsubsection{Supporting Functions}

\begin{description}
    \item[\texttt{get\_pt\_probs\_preds(model, loader, device)}]
          Returns \texttt{(true\_labels, argmax\_preds, softmax\_probs)} as
          Python lists. Used for evaluation after training.
    \item[\texttt{evaluate\_pytorch\_model(model, val\_loader, class\_names, device)}]
          Prints a full \texttt{sklearn} classification report and renders a
          confusion matrix.
    \item[\texttt{plot\_pytorch\_history(history)}]
          Renders train/validation accuracy and loss curves from the history
          dict returned by \texttt{train\_pytorch\_model}.
\end{description}

% ============================================================
\subsection{Dataset Acquisition}
\label{sec:pg-data}
% ============================================================

\subsubsection{Bug-Bite Dataset}
\label{sec:pg-data-bug}

The bug-bite dataset contains images of five species categories:
\texttt{ants}, \texttt{bed\_bugs}, \texttt{mosquitos}, \texttt{spiders},
and \texttt{ticks\_fleas}.
It is organised as \texttt{train/\{class\}/} and \texttt{val/\{class\}/}
subdirectories inside \texttt{Bug-Data/Multiclass\_Bug\_data/}.

The dataset is available via two channels:
\begin{enumerate}
    \item \textbf{Repository:} the \texttt{Bug-Data/} directory is included in
          the repository and can be used directly.
    \item \textbf{Pre-processed archive:} a curated and pre-processed version,
          including augmented samples used in the published experiments, is
          available on request from\\\texttt{ishay.yulzary@gmail.com}.
\end{enumerate}

Dataset paths are configured at the top of \texttt{miscellaneous\_code/run\_experiments.py}
(\texttt{BUG\_TRAIN}, \texttt{BUG\_VAL}) and in \texttt{scripts/shared.py}.
They default to WSL native paths (\texttt{/home/test/bug\_data/\{train,val\}}),
which must be updated for other environments.

\subsubsection{Cyclone Dataset}
\label{sec:pg-data-cyclone}

The cyclone dataset is derived from two public sources:
\begin{itemize}
    \item \textbf{HURSAT-B1} (Hurricane Satellite, Band 1): infrared brightness
          temperature satellite imagery at $8\,\text{km}$ resolution,
          hosted by NOAA.
    \item \textbf{IBTrACS} (International Best Track Archive for Climate
          Stewardship): per-observation wind speed records used to assign
          intensity labels to each satellite frame.
\end{itemize}

Five intensity categories are defined, mapped to the Saffir-Simpson scale:

\begin{center}
\begin{tabular}{lll}
\toprule
\textbf{Category label} & \textbf{Wind speed (kt)} & \textbf{Saffir-Simpson equivalent} \\
\midrule
\texttt{cat1\_tropical\_depression} & 0--33   & Tropical Depression \\
\texttt{cat2\_tropical\_storm}      & 34--63  & Tropical Storm \\
\texttt{cat3\_hurricane\_cat12}     & 64--95  & Hurricane Cat 1--2 \\
\texttt{cat4\_hurricane\_cat34}     & 96--136 & Hurricane Cat 3--4 \\
\texttt{cat5\_hurricane\_cat5}      & 137+    & Hurricane Cat 5 \\
\bottomrule
\end{tabular}
\end{center}

Images cover the years 1995--2016. Each frame is clipped to the brightness
temperature range $[180, 300]\,\text{K}$, inverted (so cold cloud tops appear
bright), and saved as a $224\times224$ PNG.

\subsubsection*{Fetching the dataset}

\begin{lstlisting}[language=bash]
# Full dataset: 4000 images per category (20000 total, ~4h on a stable connection)
python miscellaneous_code/fetch_cyclone_dataset.py

# Quick trial run: 5 images per category
python miscellaneous_code/fetch_cyclone_dataset.py --trial

# Custom target size and worker count
python miscellaneous_code/fetch_cyclone_dataset.py --target 2000 --workers 8
\end{lstlisting}

The script is resumable: interrupted runs pick up from where they stopped
using \\\texttt{cyclone\_dataset/download\_log.json}. The output directory
defaults to \texttt{cyclone\_dataset/} in the working directory.

After fetching, the dataset must be split into train/val and placed at the
paths expected by the sweep runner. The pre-built split (500 val images per
category) is available from \texttt{ishay.yulzary@gmail.com}.

% ============================================================
\subsection{Training Pipeline: Individual Scripts}
\label{sec:pg-scripts}
% ============================================================

The four scripts in \texttt{scripts/} mirror the four sections of the main
notebook. They accept the same arguments and produce the same outputs, but
are designed for unattended server execution.

\subsubsection{\texttt{01\_train\_control.py}}

Trains the three Control models (ConvNeXt-Tiny, DenseNet-121, InceptionV3)
for one seed. Saves checkpoint files at the paths supplied via CLI arguments.

\begin{lstlisting}[language=bash]
python scripts/01_train_control.py \
    --run-id s0_ls0.10_lr5.0e-06 \
    --seed 0 \
    --ctrl-convnext  results/checkpoints/control_seed_0/control_convnext.pt \
    --ctrl-densenet  results/checkpoints/control_seed_0/control_densenet.pt \
    --ctrl-inception results/checkpoints/control_seed_0/control_inception.pt
\end{lstlisting}

\subsubsection{\texttt{02\_train\_tc.py}}

For one configuration, runs cyclone pre-training on all three architectures
and then fine-tunes each cyclone backbone on the bug-bite task.
Requires both cyclone checkpoint paths (output) and TC checkpoint paths (output).

\begin{lstlisting}[language=bash]
python scripts/02_train_tc.py \
    --run-id s0_ls0.10_lr5.0e-06 \
    --seed 0 \
    --label-smoothing 0.1 \
    --phase2-lr 5e-6 \
    --cyc-convnext  results/checkpoints/config_000/cyclone_convnext.pt \
    --cyc-densenet  results/checkpoints/config_000/cyclone_densenet.pt \
    --cyc-inception results/checkpoints/config_000/cyclone_inception.pt \
    --tc-convnext   results/checkpoints/config_000/tc_bug_convnext.pt \
    --tc-densenet   results/checkpoints/config_000/tc_bug_densenet.pt \
    --tc-inception  results/checkpoints/config_000/tc_bug_inception.pt
\end{lstlisting}

\subsubsection{\texttt{03\_evaluate.py}}

Generates feature map visualisations and GradCAM overlays comparing the
Control backbone, the TC cyclone backbone (before bug-bite fine-tuning),
and the TC fine-tuned model. Outputs are saved to
\texttt{results/feature\_maps/\{run-id\}/} and
\texttt{results/plots/\{run-id\}/}.

\subsubsection{\texttt{04\_xai.py}}

Computes the full XAI suite for one configuration:
ensemble accuracy (Control vs.\ TC), silhouette scores on backbone embeddings,
t-SNE projections, LIME superpixel explanations, SHAP gradient attributions,
and DiCE counterfactual explanations. Outputs are saved to
\texttt{results/xai/\{run-id\}/}.
Optional XAI libraries (\texttt{shap}, \texttt{lime}, \texttt{dice-ml}) are
checked at runtime; cells that depend on missing libraries are skipped
gracefully.

% ============================================================
\subsection{Automated Sweep Runner: \texttt{run\_experiments.py}}
\label{sec:pg-sweep}
% ============================================================

\texttt{miscellaneous\_code/run\_experiments.py} is the primary entry point for
training. It iterates over the full 63-configuration hyperparameter grid, calls
the four pipeline scripts in sequence for each configuration, and manages
checkpointing, state, and results atomically.

\subsubsection{Basic Usage}

\begin{lstlisting}[language=bash]
# Run all pending configurations
python miscellaneous_code/run_experiments.py

# Check progress without running
python miscellaneous_code/run_experiments.py --status

# Wipe state and results (checkpoints are preserved)
python miscellaneous_code/run_experiments.py --reset
\end{lstlisting}

\subsubsection{Configuration Grid}

The grid is defined by three constants at the top of the file:

\begin{lstlisting}[language=Python]
LABEL_SMOOTHING_VALUES = [0.0, 0.1, 0.2]
PHASE2_LR_VALUES       = [1e-6, 5e-6, 1e-5]
N_SEEDS                = 7
\end{lstlisting}

Adding a value to \texttt{LABEL\_SMOOTHING\_VALUES} or \texttt{PHASE2\_LR\_VALUES}
and re-running automatically appends the new configurations to the existing
state file without invalidating completed work.

\subsubsection{Pause and Resume}

The sweep can be interrupted and resumed at any point:

\begin{description}
    \item[Graceful pause between models:] create the file
          \texttt{results/PAUSE} (\texttt{touch results/PAUSE}) or press
          Ctrl-C. The runner finishes the current model, saves its checkpoint,
          and exits cleanly.
    \item[Crash recovery:] any configuration marked \texttt{running} in the
          state file (indicating a previous interrupted run) is automatically
          reset to \texttt{pending}. Already-saved checkpoints are reused,
          so only the incomplete model needs to be retrained.
\end{description}

\subsubsection{Output Files}

\begin{description}
    \item[\texttt{results/experiment\_state.json}]
          Tracks the status of every configuration (\texttt{pending},
          \texttt{running}, or \texttt{completed}).
    \item[\texttt{results/experiment\_results.json}]
          Per-model, per-configuration evaluation results: validation accuracy,
          macro F1, precision, recall, and full classification report.
    \item[\texttt{results/checkpoints/config\_\{id\}/}]
          TC cyclone checkpoint (\texttt{cyclone\_\{model\}.pt}) and
          TC bug-bite checkpoint (\texttt{tc\_bug\_\{model\}.pt}) for each
          configuration.
    \item[\texttt{results/checkpoints/control\_seed\_\{s\}/}]
          Control checkpoints, keyed by seed only (the Control pipeline is
          independent of the cyclone hyperparameters).
    \item[\texttt{results/xai/\{run-id\}/}]
          Per-configuration XAI outputs: \texttt{ensemble\_metrics.json},
          \\\texttt{silhouette.json}, t-SNE plots, GradCAM overlays,
          LIME visualisations, SHAP attribution maps.
\end{description}

\subsubsection{Hardware Requirements}

The full 63-configuration sweep was run on an NVIDIA RTX A5000 (24\,GB VRAM).
A GPU with at least 16\,GB VRAM is recommended for training. Estimated runtime
per configuration is 1.5--2 hours on equivalent hardware.
Inference and visualisation (scripts \texttt{03} and \texttt{04}) can be run
on CPU or a smaller GPU.

% ============================================================
\subsection{Experimental Results}
\label{sec:pg-results}
% ============================================================

\subsubsection{Ensemble Accuracy: Control vs.\ TC}

The table below reports the mean ensemble validation accuracy across all seven
seeds for each hyperparameter configuration. The ensemble averages softmax
probabilities across all three architectures before taking the argmax.

\begin{center}
\begin{tabular}{lcrrrr}
\toprule
\textbf{Config (ls, lr)} & \textbf{N seeds} &
\textbf{Control mean} & \textbf{TC mean} & \textbf{$\Delta$} \\
\midrule
ls=0.10, lr=$5\times10^{-6}$ & 7 & 81.83\% & 83.27\% & +1.45 pp \\
ls=0.00, lr=$1\times10^{-6}$ & 7 & 81.83\% & 83.18\% & +1.36 pp \\
ls=0.20, lr=$1\times10^{-6}$ & 7 & 81.83\% & 83.18\% & +1.36 pp \\
ls=0.20, lr=$5\times10^{-6}$ & 7 & 81.83\% & 83.00\% & +1.18 pp \\
ls=0.20, lr=$1\times10^{-5}$ & 7 & 81.83\% & 82.82\% & +0.99 pp \\
ls=0.00, lr=$1\times10^{-5}$ & 7 & 81.83\% & 82.73\% & +0.90 pp \\
ls=0.00, lr=$5\times10^{-6}$ & 7 & 81.83\% & 82.73\% & +0.90 pp \\
ls=0.10, lr=$1\times10^{-6}$ & 7 & 81.83\% & 82.28\% & +0.45 pp \\
ls=0.10, lr=$1\times10^{-5}$ & 7 & 81.83\% & 81.83\% & +0.00 pp \\
\bottomrule
\end{tabular}
\end{center}

TC pre-training matches or exceeds the Control baseline in all nine
configurations. The best single-run result is \textbf{86.08\%} TC ensemble
accuracy (seed 4, ls=0.00, lr=$10^{-5}$), compared to 81.01\% Control for
the same seed and configuration.

\subsubsection{Backbone Embedding Quality: Silhouette Scores}

Silhouette scores are computed on backbone embeddings extracted from the
validation set (before the classification head). Higher scores indicate
better class separation in the feature space.

\begin{center}
\begin{tabular}{lrr}
\toprule
\textbf{Architecture} & \textbf{Control} & \textbf{TC} \\
\midrule
ConvNeXt-Tiny & 0.057 & 0.081 \\
DenseNet-121  & 0.027 & 0.030 \\
InceptionV3   & \text{-}0.003 & \text{-}0.017 \\
\bottomrule
\end{tabular}
\end{center}

(Values shown are from the representative seed 0, ls=0.10, lr=$5\times10^{-6}$
configuration. Per-configuration scores are available in
\texttt{results/xai/\{run-id\}/silhouette.json}.)

% ============================================================
\subsection{XAI Methods}
\label{sec:pg-xai}
% ============================================================

All XAI outputs are generated by \texttt{scripts/04\_xai.py} and by Section 4
of the main notebook.

\begin{description}
    \item[GradCAM] Class Activation Mapping via gradients. Implemented in
          \texttt{scripts/shared.py}. Target layers are architecture-specific:
          the final stage for ConvNeXt, \texttt{norm5} for DenseNet,
          and \texttt{Mixed\_7c} for InceptionV3. Produces side-by-side
          Control vs.\ TC attention overlays for direct visual comparison.

    \item[LIME] Superpixel-based local explanation using
          \texttt{lime.lime\_image.LimeImageExplainer}. The ensemble prediction
          function is used as the black-box model. Top-$N$ segments
          (default 12) are coloured by their contribution sign and weight,
          using a red--blue diverging palette. Computationally expensive at
          high \texttt{num\_samples} values; the notebook default is 1000
          samples per image.

    \item[SHAP] Pixel-level attribution via \texttt{shap.GradientExplainer}.
          A background distribution of 50 validation images is used.
          Attribution maps are averaged across the three ensemble models
          and overlaid on the input image.

    \item[DiCE] Counterfactual explanations in embedding space. Backbone
          embeddings are reduced with PCA (50 components), and a surrogate
          logistic regression classifier is fitted on the validation set.
          DiCE generates counterfactual embedding vectors showing the minimal
          directional change required to flip the predicted class.

    \item[t-SNE] 2D projection of backbone embeddings for all five bug-bite
          classes, displayed as a $3 \times 2$ grid (one row per architecture,
          columns for Control and TC). Used to visualise whether TC pre-training
          produces more separable cluster structure.

    \item[Ensemble comparison] Soft-vote ensemble across all three architectures,
          evaluated on the validation set. Control and TC ensembles are compared
          directly; results saved to \texttt{ensemble\_metrics.json}.
\end{description}

% ============================================================
\subsection{Demo Notebook: \texttt{Stacked\_Model\_Colab.ipynb}}
\label{sec:pg-demo}
% ============================================================

This notebook provides a self-contained inference and visualisation interface
designed to run on Google Colab free tier. It does not require a local GPU or
local dataset.

\paragraph{Prerequisites.}
\begin{enumerate}
    \item A Google account with access to the project Google Drive folder.
          The folder must contain \texttt{project\_book\_results.tar},
          which packages all trained weight files.
    \item The notebook set \texttt{PROJECT\_ROOT} to the correct Drive path
          \\(default: \texttt{/content/drive/MyDrive/Capstone}).
\end{enumerate}

\paragraph{Workflow.}
\begin{enumerate}
    \item Mount Google Drive (Cell 2).
    \item Extract weights from the archive if not already extracted (Cell 3).
    \item Load all six models: three Control and three TC (Cell 5).
    \item Set \texttt{IMG\_PATH} to any JPEG or PNG image (Cell 7).
    \item Run inference: Cell 7 displays prediction probabilities for both
          pipelines side-by-side.
    \item Run GradCAM: Cell 8 renders a $3 \times 2$ overlay grid
          (architectures $\times$ pipelines).
    \item Run LIME: Cell 9 renders superpixel attributions for both pipelines.
\end{enumerate}

Weight file locations within the extracted archive:
\begin{itemize}
    \item Control: \texttt{project\_book\_results/s3\_ls0.20\_lr1.0e-05/}
    \item TC: \texttt{project\_book\_results/s6\_ls0.10\_lr5.0e-06/}
\end{itemize}

% ============================================================
\subsection{Extending the System}
\label{sec:pg-extend}
% ============================================================

\subsubsection{Adding a New Architecture}

\begin{enumerate}
    \item Add an entry to the \texttt{MODELS} list in
          \texttt{miscellaneous\_code/run\_experiments.py}:
          \texttt{(key, timm\_id, img\_size, cyc\_p1\_bs, bug\_bs,
          cyc\_p2\_epochs, bug\_p2\_epochs)}.
    \item Identify the correct GradCAM target layer for the new architecture
          and add a case to the \texttt{\_get\_target} function in
          \texttt{Stacked\_Model\_Colab.ipynb} and in \texttt{scripts/shared.py}.
    \item Re-run the sweep. New configurations are appended automatically;
          existing results are not invalidated.
\end{enumerate}

\subsubsection{Adding a New Bug-Bite Class}

\begin{enumerate}
    \item Add labelled images to \texttt{Bug-Data/Multiclass\_Bug\_data/train/}
          and \texttt{val/} under a new subdirectory named after the class.
    \item Update \texttt{N\_CLASSES} in \texttt{run\_experiments.py} if
          necessary (currently 5; \texttt{timm} models are instantiated
          with this value).
    \item Retrain from scratch. The existing Control and TC checkpoints are
          incompatible with a different number of output classes.
\end{enumerate}

\subsubsection{Adapting to a New Intermediate Domain}

The TC pipeline is not specific to cyclone imagery. Any image classification
task can serve as the intermediate pre-training domain by:
\begin{enumerate}
    \item Placing the new domain dataset at \texttt{TC\_TRAIN} / \texttt{TC\_VAL}
          in \texttt{run\_experiments.py}.
    \item Adjusting the augmentation strategy in \texttt{get\_pytorch\_loaders}
          if the new domain does not share cyclone imagery's symmetry properties.
    \item Updating the \texttt{augment='strong'} branch of
          \texttt{pytorch\_utils.py} accordingly.
\end{enumerate}

\subsubsection{Path Toward Productisation}

The current system is a research prototype. The following steps are required
before deployment:

\begin{enumerate}
    \item \textbf{Model serving:} wrap the ensemble inference logic from
          \texttt{Stacked\_Model\_Colab.ipynb} in a REST API endpoint
          (e.g.\ FastAPI). The core inference function is approximately
          30 lines of PyTorch code.
    \item \textbf{Weight hosting:} move trained weights from Google Drive to
          a persistent object store (e.g.\ AWS S3, GCP Cloud Storage).
    \item \textbf{Input validation:} add image size checks, colour space
          normalisation, and rejection of non-photographic inputs.
    \item \textbf{Dataset expansion:} the current bug-bite dataset is small
          (fewer than 500 images per class). Classification performance will
          improve substantially with a larger, clinically curated dataset.
    \item \textbf{Clinical validation:} model outputs should be reviewed and
          validated by a medical professional before any medical application.
\end{enumerate}

% --- End of appendix (remove if standalone) ---
\end{document}

%
%  Or compile standalone:
%    Uncomment the preamble block below, and add \end{document} at the end.
% ============================================================

% --- Standalone preamble (remove if including in main document) ----------
\documentclass[12pt,a4paper]{article}
\usepackage[utf8]{inputenc}
\usepackage[T1]{fontenc}
\usepackage{geometry}
\geometry{margin=2.5cm}
\usepackage{hyperref}
\usepackage{listings}
\usepackage{xcolor}
\usepackage{booktabs}
\usepackage{array}
\usepackage{enumitem}
\usepackage{parskip}
\lstset{basicstyle=\ttfamily\small, breaklines=true, frame=single,
        backgroundcolor=\color{gray!8}}
\begin{document}
\title{Programmer's Guide\\
       \large Bug-Bite Classification via Tropical Cyclone Transfer Learning\\
       Braude College of Engineering --- Software Engineering Capstone}
\author{Ishay Yulzary, Nathan Trostianitser}
\date{}
\maketitle
\tableofcontents
\newpage
% -------------------------------------------------------------------------

\section{Programmer's Guide}
\label{app:programmer-guide}

This appendix is addressed to a software engineer who wishes to understand the
internals of the system, reproduce the experimental results, extend the research,
or advance the prototype toward a deployable product.
All source code lives in the repository root unless otherwise noted.

% ============================================================
\subsection{Research Overview}
\label{sec:pg-overview}
% ============================================================

\subsubsection{Problem and Motivation}

Classifying bug-bite skin images by causative species is clinically useful but
technically difficult: inter-class visual similarity is high, intra-class variation
is large, and labelled medical image datasets are small. Standard fine-tuning from
ImageNet leaves the backbone biased toward object-centric features that are of
limited relevance to dermatological macro-photography.

\subsubsection{Core Hypothesis}

This project investigates whether inserting an intermediate pre-training step on
tropical cyclone (TC) infrared satellite imagery improves backbone representations
for the bug-bite task. The conjecture is that cyclone satellite images and skin
macro-photographs share structural visual properties: radial textures, soft
gradients, and fine-grained pattern differences between severity categories that
are invisible to coarse object detectors. A backbone that has already learned to
discriminate these properties at the cyclone level may require less adaptation to
do so at the skin level.

\subsubsection{Transfer Classification Pipeline}

Two training pipelines are compared:

\begin{description}
    \item[Control] ImageNet pre-training $\rightarrow$ bug-bite fine-tuning
          (direct, single-stage transfer).
    \item[TC] ImageNet pre-training $\rightarrow$ cyclone intensity pre-training
          $\rightarrow$ bug-bite fine-tuning (two-stage transfer).
\end{description}

The cyclone pre-training stage is a full classification task: the model learns to
distinguish five cyclone intensity categories derived from the Saffir-Simpson
scale. The resulting backbone weights are then transplanted into a fresh
bug-bite classification head, and the whole network is fine-tuned.

\subsubsection{Experimental Design}

Three backbone architectures are evaluated under both pipelines:

\begin{center}
\begin{tabular}{lll}
\toprule
\textbf{Architecture} & \textbf{timm identifier} & \textbf{Input size} \\
\midrule
ConvNeXt-Tiny  & \texttt{convnext\_tiny.fb\_in22k\_ft\_in1k} & 256$\times$256 \\
DenseNet-121   & \texttt{densenet121.ra\_in1k}                & 256$\times$256 \\
InceptionV3    & \texttt{inception\_v3.tv\_in1k}              & 299$\times$299 \\
\bottomrule
\end{tabular}
\end{center}

The cyclone pre-training hyperparameters are swept across the following grid,
yielding 63 configurations in total (3 label-smoothing values $\times$ 3 learning
rates $\times$ 7 random seeds):

\begin{center}
\begin{tabular}{ll}
\toprule
\textbf{Parameter} & \textbf{Values} \\
\midrule
Cyclone label smoothing (\texttt{ls}) & 0.0, 0.1, 0.2 \\
Cyclone phase-2 learning rate (\texttt{lr}) & $10^{-6}$, $5\times10^{-6}$, $10^{-5}$ \\
Random seeds & 0, 1, 2, 3, 4, 5, 6 \\
\bottomrule
\end{tabular}
\end{center}

Bug-bite fine-tuning hyperparameters are fixed across all configurations:
label smoothing 0.0, phase-2 learning rate $5\times10^{-6}$.

% ============================================================
\subsection{Repository Structure}
\label{sec:pg-repo}
% ============================================================

\begin{description}
    \item[\texttt{notebooks/Multiclass\_Classification.ipynb}]
          The primary experimental notebook. Covers all four pipeline stages:
          control training, TC pre-training, feature map and GradCAM evaluation,
          and the full XAI suite. Intended for interactive exploration or
          one-shot reproduction on a GPU instance.

    \item[\texttt{Stacked\_Model\_Colab.ipynb}]
          Self-contained inference and visualisation notebook designed for
          Google Colab free tier. Loads trained weights from Google Drive,
          runs ensemble inference on a single image, and produces side-by-side
          Control vs.\ TC GradCAM and LIME outputs.

    \item[\texttt{scripts/}]
          Standalone Python scripts, one per pipeline stage. Intended for
          execution on a headless server via \texttt{run\_experiments.py}.
          \begin{itemize}
              \item \texttt{shared.py} --- shared imports, dataset paths,
                    metrics helpers, and GradCAM utilities used by all four scripts.
              \item \texttt{01\_train\_control.py} --- trains the Control
                    (ImageNet $\rightarrow$ bug-bite) models.
              \item \texttt{02\_train\_tc.py} --- runs cyclone pre-training then
                    TC $\rightarrow$ bug-bite fine-tuning.
              \item \texttt{03\_evaluate.py} --- generates feature map comparisons
                    and GradCAM overlays.
              \item \texttt{04\_xai.py} --- computes ensemble metrics, silhouette
                    scores, t-SNE projections, LIME, SHAP, and DiCE outputs.
          \end{itemize}

    \item[\texttt{miscellaneous\_code/pytorch\_utils.py}]
          Custom PyTorch training library imported by both the notebook and the
          scripts. See Section~\ref{sec:pg-pytorch-utils} for full documentation.

    \item[\texttt{miscellaneous\_code/run\_experiments.py}]
          Automated hyperparameter sweep runner. See
          Section~\ref{sec:pg-sweep} for full documentation.

    \item[\texttt{miscellaneous\_code/fetch\_cyclone\_dataset.py}]
          Downloads and labels the HURSAT-B1 cyclone dataset from NOAA.
          See Section~\ref{sec:pg-data-cyclone}.

    \item[\texttt{miscellaneous\_code/cyclone\_preprocessing.py}]
          Removes grid-line artefacts introduced by the HURSAT-B1 satellite
          image tiling process before training.

    \item[\texttt{Bug-Data/Multiclass\_Bug\_data/}]
          Bug-bite dataset, structured as \texttt{train/} and \texttt{val/}
          subdirectories, each containing one subdirectory per class.

    \item[\texttt{results/}]
          Auto-generated by the sweep runner. Contains checkpoints, per-run
          metric JSONs, XAI outputs, and plots.

    \item[\texttt{deploy/}]
          Legacy Streamlit demo (binary classifier, outdated). Not part of
          the active pipeline; retained for reference only.
\end{description}

% ============================================================
\subsection{Core Training Library: \texttt{pytorch\_utils.py}}
\label{sec:pg-pytorch-utils}
% ============================================================

This file contains all reusable PyTorch infrastructure. It is imported by
\texttt{run\_experiments.py}, by the four pipeline scripts in \texttt{scripts/},
and by the main notebook. Engineers extending the system should understand each
public function before modifying the training loop.

\subsubsection{\texttt{get\_pytorch\_loaders}}

\begin{lstlisting}[language=Python]
get_pytorch_loaders(train_dir, val_dir,
                    img_size=310, batch_size=4, augment=False)
\end{lstlisting}

Returns \texttt{(train\_loader, val\_loader, class\_names)} using
\texttt{torchvision.datasets.ImageFolder}. The \texttt{augment} parameter
accepts three values:

\begin{description}
    \item[\texttt{False}] No augmentation. Used for validation and for
          bug-bite fine-tuning in the sweep (the bug-bite dataset is small;
          augmentation during fine-tuning degraded performance in preliminary
          experiments).
    \item[\texttt{True}] Light augmentation: horizontal flip, $\pm10^{\circ}$
          rotation. Used for moderate augmentation scenarios.
    \item[\texttt{'strong'}] Aggressive augmentation designed for cyclone
          satellite imagery: random resized crop (scale 0.35--1.0), full
          $360^{\circ}$ rotation, vertical flip, colour jitter, random greyscale,
          Gaussian blur (kernel 5, $\sigma \in [0.5, 3.0]$), and random erasing.
          The full rotation and vertical flip are valid for cyclone images because
          satellite imagery is radially symmetric with respect to the storm centre.
\end{description}

All loaders use ImageNet normalisation
(\texttt{mean=[0.485, 0.456, 0.406]}, \texttt{std=[0.229, 0.224, 0.225]}).

\subsubsection{\texttt{train\_pytorch\_model}}

\begin{lstlisting}[language=Python]
train_pytorch_model(model, train_loader, val_loader, device,
                    phase1_epochs=5, phase2_epochs=20, patience=5,
                    save_path=None,
                    phase1_batch_size=None, phase2_batch_size=None,
                    label_smoothing=0.0, auto_class_weights=False,
                    phase1_lr=1e-4, phase2_lr=5e-6)
\end{lstlisting}

Implements the standard two-phase fine-tuning strategy used throughout the
project:

\begin{description}
    \item[Phase 1 (head-only)] All backbone parameters are frozen.
          Only the classification head (detected via \texttt{model.get\_classifier()})
          is trained for \texttt{phase1\_epochs} epochs at \texttt{phase1\_lr}
          with AdamW ($\lambda = 0.01$).
    \item[Phase 2 (full fine-tune)] All parameters are unfrozen and the
          entire network is trained at \texttt{phase2\_lr}. Gradient checkpointing
          is enabled where the model supports it.
\end{description}

Both phases use:
\begin{itemize}
    \item Early stopping on validation loss with patience \texttt{patience}.
    \item Automatic Mixed Precision (AMP) with \texttt{bfloat16} when a CUDA
          device is available.
    \item \texttt{CrossEntropyLoss} with optional label smoothing.
    \item Patience counters that reset independently between phases,
          because phase 2 typically causes an initial loss spike after
          unfreezing the backbone.
\end{itemize}

The best checkpoint (lowest validation loss) is saved to \texttt{save\_path}
and automatically reloaded into the model before the function returns.

\subsubsection{\texttt{get\_backbone\_state\_dict}}

\begin{lstlisting}[language=Python]
get_backbone_state_dict(state_dict)
\end{lstlisting}

Strips all classifier head keys from a state dict, enabling backbone weight
transfer between tasks with different numbers of output classes. Keys whose
names start with \texttt{head.}, \texttt{fc.}, or \texttt{classifier.} are
removed. This function is the mechanism by which the cyclone-pre-trained
backbone is transplanted into a fresh bug-bite classification head.

\subsubsection{Supporting Functions}

\begin{description}
    \item[\texttt{get\_pt\_probs\_preds(model, loader, device)}]
          Returns \texttt{(true\_labels, argmax\_preds, softmax\_probs)} as
          Python lists. Used for evaluation after training.
    \item[\texttt{evaluate\_pytorch\_model(model, val\_loader, class\_names, device)}]
          Prints a full \texttt{sklearn} classification report and renders a
          confusion matrix.
    \item[\texttt{plot\_pytorch\_history(history)}]
          Renders train/validation accuracy and loss curves from the history
          dict returned by \texttt{train\_pytorch\_model}.
\end{description}

% ============================================================
\subsection{Dataset Acquisition}
\label{sec:pg-data}
% ============================================================

\subsubsection{Bug-Bite Dataset}
\label{sec:pg-data-bug}

The bug-bite dataset contains images of five species categories:
\texttt{ants}, \texttt{bed\_bugs}, \texttt{mosquitos}, \texttt{spiders},
and \texttt{ticks\_fleas}.
It is organised as \texttt{train/\{class\}/} and \texttt{val/\{class\}/}
subdirectories inside \texttt{Bug-Data/Multiclass\_Bug\_data/}.

The dataset is available via two channels:
\begin{enumerate}
    \item \textbf{Repository:} the \texttt{Bug-Data/} directory is included in
          the repository and can be used directly.
    \item \textbf{Pre-processed archive:} a curated and pre-processed version,
          including augmented samples used in the published experiments, is
          available on request from\\\texttt{ishay.yulzary@gmail.com}.
\end{enumerate}

Dataset paths are configured at the top of \texttt{miscellaneous\_code/run\_experiments.py}
(\texttt{BUG\_TRAIN}, \texttt{BUG\_VAL}) and in \texttt{scripts/shared.py}.
They default to WSL native paths (\texttt{/home/test/bug\_data/\{train,val\}}),
which must be updated for other environments.

\subsubsection{Cyclone Dataset}
\label{sec:pg-data-cyclone}

The cyclone dataset is derived from two public sources:
\begin{itemize}
    \item \textbf{HURSAT-B1} (Hurricane Satellite, Band 1): infrared brightness
          temperature satellite imagery at $8\,\text{km}$ resolution,
          hosted by NOAA.
    \item \textbf{IBTrACS} (International Best Track Archive for Climate
          Stewardship): per-observation wind speed records used to assign
          intensity labels to each satellite frame.
\end{itemize}

Five intensity categories are defined, mapped to the Saffir-Simpson scale:

\begin{center}
\begin{tabular}{lll}
\toprule
\textbf{Category label} & \textbf{Wind speed (kt)} & \textbf{Saffir-Simpson equivalent} \\
\midrule
\texttt{cat1\_tropical\_depression} & 0--33   & Tropical Depression \\
\texttt{cat2\_tropical\_storm}      & 34--63  & Tropical Storm \\
\texttt{cat3\_hurricane\_cat12}     & 64--95  & Hurricane Cat 1--2 \\
\texttt{cat4\_hurricane\_cat34}     & 96--136 & Hurricane Cat 3--4 \\
\texttt{cat5\_hurricane\_cat5}      & 137+    & Hurricane Cat 5 \\
\bottomrule
\end{tabular}
\end{center}

Images cover the years 1995--2016. Each frame is clipped to the brightness
temperature range $[180, 300]\,\text{K}$, inverted (so cold cloud tops appear
bright), and saved as a $224\times224$ PNG.

\subsubsection*{Fetching the dataset}

\begin{lstlisting}[language=bash]
# Full dataset: 4000 images per category (20000 total, ~4h on a stable connection)
python miscellaneous_code/fetch_cyclone_dataset.py

# Quick trial run: 5 images per category
python miscellaneous_code/fetch_cyclone_dataset.py --trial

# Custom target size and worker count
python miscellaneous_code/fetch_cyclone_dataset.py --target 2000 --workers 8
\end{lstlisting}

The script is resumable: interrupted runs pick up from where they stopped
using \\\texttt{cyclone\_dataset/download\_log.json}. The output directory
defaults to \texttt{cyclone\_dataset/} in the working directory.

After fetching, the dataset must be split into train/val and placed at the
paths expected by the sweep runner. The pre-built split (500 val images per
category) is available from \texttt{ishay.yulzary@gmail.com}.

% ============================================================
\subsection{Training Pipeline: Individual Scripts}
\label{sec:pg-scripts}
% ============================================================

The four scripts in \texttt{scripts/} mirror the four sections of the main
notebook. They accept the same arguments and produce the same outputs, but
are designed for unattended server execution.

\subsubsection{\texttt{01\_train\_control.py}}

Trains the three Control models (ConvNeXt-Tiny, DenseNet-121, InceptionV3)
for one seed. Saves checkpoint files at the paths supplied via CLI arguments.

\begin{lstlisting}[language=bash]
python scripts/01_train_control.py \
    --run-id s0_ls0.10_lr5.0e-06 \
    --seed 0 \
    --ctrl-convnext  results/checkpoints/control_seed_0/control_convnext.pt \
    --ctrl-densenet  results/checkpoints/control_seed_0/control_densenet.pt \
    --ctrl-inception results/checkpoints/control_seed_0/control_inception.pt
\end{lstlisting}

\subsubsection{\texttt{02\_train\_tc.py}}

For one configuration, runs cyclone pre-training on all three architectures
and then fine-tunes each cyclone backbone on the bug-bite task.
Requires both cyclone checkpoint paths (output) and TC checkpoint paths (output).

\begin{lstlisting}[language=bash]
python scripts/02_train_tc.py \
    --run-id s0_ls0.10_lr5.0e-06 \
    --seed 0 \
    --label-smoothing 0.1 \
    --phase2-lr 5e-6 \
    --cyc-convnext  results/checkpoints/config_000/cyclone_convnext.pt \
    --cyc-densenet  results/checkpoints/config_000/cyclone_densenet.pt \
    --cyc-inception results/checkpoints/config_000/cyclone_inception.pt \
    --tc-convnext   results/checkpoints/config_000/tc_bug_convnext.pt \
    --tc-densenet   results/checkpoints/config_000/tc_bug_densenet.pt \
    --tc-inception  results/checkpoints/config_000/tc_bug_inception.pt
\end{lstlisting}

\subsubsection{\texttt{03\_evaluate.py}}

Generates feature map visualisations and GradCAM overlays comparing the
Control backbone, the TC cyclone backbone (before bug-bite fine-tuning),
and the TC fine-tuned model. Outputs are saved to
\texttt{results/feature\_maps/\{run-id\}/} and
\texttt{results/plots/\{run-id\}/}.

\subsubsection{\texttt{04\_xai.py}}

Computes the full XAI suite for one configuration:
ensemble accuracy (Control vs.\ TC), silhouette scores on backbone embeddings,
t-SNE projections, LIME superpixel explanations, SHAP gradient attributions,
and DiCE counterfactual explanations. Outputs are saved to
\texttt{results/xai/\{run-id\}/}.
Optional XAI libraries (\texttt{shap}, \texttt{lime}, \texttt{dice-ml}) are
checked at runtime; cells that depend on missing libraries are skipped
gracefully.

% ============================================================
\subsection{Automated Sweep Runner: \texttt{run\_experiments.py}}
\label{sec:pg-sweep}
% ============================================================

\texttt{miscellaneous\_code/run\_experiments.py} is the primary entry point for
training. It iterates over the full 63-configuration hyperparameter grid, calls
the four pipeline scripts in sequence for each configuration, and manages
checkpointing, state, and results atomically.

\subsubsection{Basic Usage}

\begin{lstlisting}[language=bash]
# Run all pending configurations
python miscellaneous_code/run_experiments.py

# Check progress without running
python miscellaneous_code/run_experiments.py --status

# Wipe state and results (checkpoints are preserved)
python miscellaneous_code/run_experiments.py --reset
\end{lstlisting}

\subsubsection{Configuration Grid}

The grid is defined by three constants at the top of the file:

\begin{lstlisting}[language=Python]
LABEL_SMOOTHING_VALUES = [0.0, 0.1, 0.2]
PHASE2_LR_VALUES       = [1e-6, 5e-6, 1e-5]
N_SEEDS                = 7
\end{lstlisting}

Adding a value to \texttt{LABEL\_SMOOTHING\_VALUES} or \texttt{PHASE2\_LR\_VALUES}
and re-running automatically appends the new configurations to the existing
state file without invalidating completed work.

\subsubsection{Pause and Resume}

The sweep can be interrupted and resumed at any point:

\begin{description}
    \item[Graceful pause between models:] create the file
          \texttt{results/PAUSE} (\texttt{touch results/PAUSE}) or press
          Ctrl-C. The runner finishes the current model, saves its checkpoint,
          and exits cleanly.
    \item[Crash recovery:] any configuration marked \texttt{running} in the
          state file (indicating a previous interrupted run) is automatically
          reset to \texttt{pending}. Already-saved checkpoints are reused,
          so only the incomplete model needs to be retrained.
\end{description}

\subsubsection{Output Files}

\begin{description}
    \item[\texttt{results/experiment\_state.json}]
          Tracks the status of every configuration (\texttt{pending},
          \texttt{running}, or \texttt{completed}).
    \item[\texttt{results/experiment\_results.json}]
          Per-model, per-configuration evaluation results: validation accuracy,
          macro F1, precision, recall, and full classification report.
    \item[\texttt{results/checkpoints/config\_\{id\}/}]
          TC cyclone checkpoint (\texttt{cyclone\_\{model\}.pt}) and
          TC bug-bite checkpoint (\texttt{tc\_bug\_\{model\}.pt}) for each
          configuration.
    \item[\texttt{results/checkpoints/control\_seed\_\{s\}/}]
          Control checkpoints, keyed by seed only (the Control pipeline is
          independent of the cyclone hyperparameters).
    \item[\texttt{results/xai/\{run-id\}/}]
          Per-configuration XAI outputs: \texttt{ensemble\_metrics.json},
          \\\texttt{silhouette.json}, t-SNE plots, GradCAM overlays,
          LIME visualisations, SHAP attribution maps.
\end{description}

\subsubsection{Hardware Requirements}

The full 63-configuration sweep was run on an NVIDIA RTX A5000 (24\,GB VRAM).
A GPU with at least 16\,GB VRAM is recommended for training. Estimated runtime
per configuration is 1.5--2 hours on equivalent hardware.
Inference and visualisation (scripts \texttt{03} and \texttt{04}) can be run
on CPU or a smaller GPU.

% ============================================================
\subsection{Experimental Results}
\label{sec:pg-results}
% ============================================================

\subsubsection{Ensemble Accuracy: Control vs.\ TC}

The table below reports the mean ensemble validation accuracy across all seven
seeds for each hyperparameter configuration. The ensemble averages softmax
probabilities across all three architectures before taking the argmax.

\begin{center}
\begin{tabular}{lcrrrr}
\toprule
\textbf{Config (ls, lr)} & \textbf{N seeds} &
\textbf{Control mean} & \textbf{TC mean} & \textbf{$\Delta$} \\
\midrule
ls=0.10, lr=$5\times10^{-6}$ & 7 & 81.83\% & 83.27\% & +1.45 pp \\
ls=0.00, lr=$1\times10^{-6}$ & 7 & 81.83\% & 83.18\% & +1.36 pp \\
ls=0.20, lr=$1\times10^{-6}$ & 7 & 81.83\% & 83.18\% & +1.36 pp \\
ls=0.20, lr=$5\times10^{-6}$ & 7 & 81.83\% & 83.00\% & +1.18 pp \\
ls=0.20, lr=$1\times10^{-5}$ & 7 & 81.83\% & 82.82\% & +0.99 pp \\
ls=0.00, lr=$1\times10^{-5}$ & 7 & 81.83\% & 82.73\% & +0.90 pp \\
ls=0.00, lr=$5\times10^{-6}$ & 7 & 81.83\% & 82.73\% & +0.90 pp \\
ls=0.10, lr=$1\times10^{-6}$ & 7 & 81.83\% & 82.28\% & +0.45 pp \\
ls=0.10, lr=$1\times10^{-5}$ & 7 & 81.83\% & 81.83\% & +0.00 pp \\
\bottomrule
\end{tabular}
\end{center}

TC pre-training matches or exceeds the Control baseline in all nine
configurations. The best single-run result is \textbf{86.08\%} TC ensemble
accuracy (seed 4, ls=0.00, lr=$10^{-5}$), compared to 81.01\% Control for
the same seed and configuration.

\subsubsection{Backbone Embedding Quality: Silhouette Scores}

Silhouette scores are computed on backbone embeddings extracted from the
validation set (before the classification head). Higher scores indicate
better class separation in the feature space.

\begin{center}
\begin{tabular}{lrr}
\toprule
\textbf{Architecture} & \textbf{Control} & \textbf{TC} \\
\midrule
ConvNeXt-Tiny & 0.057 & 0.081 \\
DenseNet-121  & 0.027 & 0.030 \\
InceptionV3   & \text{-}0.003 & \text{-}0.017 \\
\bottomrule
\end{tabular}
\end{center}

(Values shown are from the representative seed 0, ls=0.10, lr=$5\times10^{-6}$
configuration. Per-configuration scores are available in
\texttt{results/xai/\{run-id\}/silhouette.json}.)

% ============================================================
\subsection{XAI Methods}
\label{sec:pg-xai}
% ============================================================

All XAI outputs are generated by \texttt{scripts/04\_xai.py} and by Section 4
of the main notebook.

\begin{description}
    \item[GradCAM] Class Activation Mapping via gradients. Implemented in
          \texttt{scripts/shared.py}. Target layers are architecture-specific:
          the final stage for ConvNeXt, \texttt{norm5} for DenseNet,
          and \texttt{Mixed\_7c} for InceptionV3. Produces side-by-side
          Control vs.\ TC attention overlays for direct visual comparison.

    \item[LIME] Superpixel-based local explanation using
          \texttt{lime.lime\_image.LimeImageExplainer}. The ensemble prediction
          function is used as the black-box model. Top-$N$ segments
          (default 12) are coloured by their contribution sign and weight,
          using a red--blue diverging palette. Computationally expensive at
          high \texttt{num\_samples} values; the notebook default is 1000
          samples per image.

    \item[SHAP] Pixel-level attribution via \texttt{shap.GradientExplainer}.
          A background distribution of 50 validation images is used.
          Attribution maps are averaged across the three ensemble models
          and overlaid on the input image.

    \item[DiCE] Counterfactual explanations in embedding space. Backbone
          embeddings are reduced with PCA (50 components), and a surrogate
          logistic regression classifier is fitted on the validation set.
          DiCE generates counterfactual embedding vectors showing the minimal
          directional change required to flip the predicted class.

    \item[t-SNE] 2D projection of backbone embeddings for all five bug-bite
          classes, displayed as a $3 \times 2$ grid (one row per architecture,
          columns for Control and TC). Used to visualise whether TC pre-training
          produces more separable cluster structure.

    \item[Ensemble comparison] Soft-vote ensemble across all three architectures,
          evaluated on the validation set. Control and TC ensembles are compared
          directly; results saved to \texttt{ensemble\_metrics.json}.
\end{description}

% ============================================================
\subsection{Demo Notebook: \texttt{Stacked\_Model\_Colab.ipynb}}
\label{sec:pg-demo}
% ============================================================

This notebook provides a self-contained inference and visualisation interface
designed to run on Google Colab free tier. It does not require a local GPU or
local dataset.

\paragraph{Prerequisites.}
\begin{enumerate}
    \item A Google account with access to the project Google Drive folder.
          The folder must contain \texttt{project\_book\_results.tar},
          which packages all trained weight files.
    \item The notebook set \texttt{PROJECT\_ROOT} to the correct Drive path
          \\(default: \texttt{/content/drive/MyDrive/Capstone}).
\end{enumerate}

\paragraph{Workflow.}
\begin{enumerate}
    \item Mount Google Drive (Cell 2).
    \item Extract weights from the archive if not already extracted (Cell 3).
    \item Load all six models: three Control and three TC (Cell 5).
    \item Set \texttt{IMG\_PATH} to any JPEG or PNG image (Cell 7).
    \item Run inference: Cell 7 displays prediction probabilities for both
          pipelines side-by-side.
    \item Run GradCAM: Cell 8 renders a $3 \times 2$ overlay grid
          (architectures $\times$ pipelines).
    \item Run LIME: Cell 9 renders superpixel attributions for both pipelines.
\end{enumerate}

Weight file locations within the extracted archive:
\begin{itemize}
    \item Control: \texttt{project\_book\_results/s3\_ls0.20\_lr1.0e-05/}
    \item TC: \texttt{project\_book\_results/s6\_ls0.10\_lr5.0e-06/}
\end{itemize}

% ============================================================
\subsection{Extending the System}
\label{sec:pg-extend}
% ============================================================

\subsubsection{Adding a New Architecture}

\begin{enumerate}
    \item Add an entry to the \texttt{MODELS} list in
          \texttt{miscellaneous\_code/run\_experiments.py}:
          \texttt{(key, timm\_id, img\_size, cyc\_p1\_bs, bug\_bs,
          cyc\_p2\_epochs, bug\_p2\_epochs)}.
    \item Identify the correct GradCAM target layer for the new architecture
          and add a case to the \texttt{\_get\_target} function in
          \texttt{Stacked\_Model\_Colab.ipynb} and in \texttt{scripts/shared.py}.
    \item Re-run the sweep. New configurations are appended automatically;
          existing results are not invalidated.
\end{enumerate}

\subsubsection{Adding a New Bug-Bite Class}

\begin{enumerate}
    \item Add labelled images to \texttt{Bug-Data/Multiclass\_Bug\_data/train/}
          and \texttt{val/} under a new subdirectory named after the class.
    \item Update \texttt{N\_CLASSES} in \texttt{run\_experiments.py} if
          necessary (currently 5; \texttt{timm} models are instantiated
          with this value).
    \item Retrain from scratch. The existing Control and TC checkpoints are
          incompatible with a different number of output classes.
\end{enumerate}

\subsubsection{Adapting to a New Intermediate Domain}

The TC pipeline is not specific to cyclone imagery. Any image classification
task can serve as the intermediate pre-training domain by:
\begin{enumerate}
    \item Placing the new domain dataset at \texttt{TC\_TRAIN} / \texttt{TC\_VAL}
          in \texttt{run\_experiments.py}.
    \item Adjusting the augmentation strategy in \texttt{get\_pytorch\_loaders}
          if the new domain does not share cyclone imagery's symmetry properties.
    \item Updating the \texttt{augment='strong'} branch of
          \texttt{pytorch\_utils.py} accordingly.
\end{enumerate}

\subsubsection{Path Toward Productisation}

The current system is a research prototype. The following steps are required
before deployment:

\begin{enumerate}
    \item \textbf{Model serving:} wrap the ensemble inference logic from
          \texttt{Stacked\_Model\_Colab.ipynb} in a REST API endpoint
          (e.g.\ FastAPI). The core inference function is approximately
          30 lines of PyTorch code.
    \item \textbf{Weight hosting:} move trained weights from Google Drive to
          a persistent object store (e.g.\ AWS S3, GCP Cloud Storage).
    \item \textbf{Input validation:} add image size checks, colour space
          normalisation, and rejection of non-photographic inputs.
    \item \textbf{Dataset expansion:} the current bug-bite dataset is small
          (fewer than 500 images per class). Classification performance will
          improve substantially with a larger, clinically curated dataset.
    \item \textbf{Clinical validation:} model outputs should be reviewed and
          validated by a medical professional before any medical application.
\end{enumerate}

% --- End of appendix (remove if standalone) ---
\end{document}

%
%  Or compile standalone:
%    Uncomment the preamble block below, and add \end{document} at the end.
% ============================================================

% --- Standalone preamble (remove if including in main document) ----------
\documentclass[12pt,a4paper]{article}
\usepackage[utf8]{inputenc}
\usepackage[T1]{fontenc}
\usepackage{geometry}
\geometry{margin=2.5cm}
\usepackage{hyperref}
\usepackage{listings}
\usepackage{xcolor}
\usepackage{booktabs}
\usepackage{array}
\usepackage{enumitem}
\usepackage{parskip}
\lstset{basicstyle=\ttfamily\small, breaklines=true, frame=single,
        backgroundcolor=\color{gray!8}}
\begin{document}
\title{Programmer's Guide\\
       \large Bug-Bite Classification via Tropical Cyclone Transfer Learning\\
       Braude College of Engineering --- Software Engineering Capstone}
\author{Ishay Yulzary, Nathan Trostianitser}
\date{}
\maketitle
\tableofcontents
\newpage
% -------------------------------------------------------------------------

\section{Programmer's Guide}
\label{app:programmer-guide}

This appendix is addressed to a software engineer who wishes to understand the
internals of the system, reproduce the experimental results, extend the research,
or advance the prototype toward a deployable product.
All source code lives in the repository root unless otherwise noted.

% ============================================================
\subsection{Research Overview}
\label{sec:pg-overview}
% ============================================================

\subsubsection{Problem and Motivation}

Classifying bug-bite skin images by causative species is clinically useful but
technically difficult: inter-class visual similarity is high, intra-class variation
is large, and labelled medical image datasets are small. Standard fine-tuning from
ImageNet leaves the backbone biased toward object-centric features that are of
limited relevance to dermatological macro-photography.

\subsubsection{Core Hypothesis}

This project investigates whether inserting an intermediate pre-training step on
tropical cyclone (TC) infrared satellite imagery improves backbone representations
for the bug-bite task. The conjecture is that cyclone satellite images and skin
macro-photographs share structural visual properties: radial textures, soft
gradients, and fine-grained pattern differences between severity categories that
are invisible to coarse object detectors. A backbone that has already learned to
discriminate these properties at the cyclone level may require less adaptation to
do so at the skin level.

\subsubsection{Transfer Classification Pipeline}

Two training pipelines are compared:

\begin{description}
    \item[Control] ImageNet pre-training $\rightarrow$ bug-bite fine-tuning
          (direct, single-stage transfer).
    \item[TC] ImageNet pre-training $\rightarrow$ cyclone intensity pre-training
          $\rightarrow$ bug-bite fine-tuning (two-stage transfer).
\end{description}

The cyclone pre-training stage is a full classification task: the model learns to
distinguish five cyclone intensity categories derived from the Saffir-Simpson
scale. The resulting backbone weights are then transplanted into a fresh
bug-bite classification head, and the whole network is fine-tuned.

\subsubsection{Experimental Design}

Three backbone architectures are evaluated under both pipelines:

\begin{center}
\begin{tabular}{lll}
\toprule
\textbf{Architecture} & \textbf{timm identifier} & \textbf{Input size} \\
\midrule
ConvNeXt-Tiny  & \texttt{convnext\_tiny.fb\_in22k\_ft\_in1k} & 256$\times$256 \\
DenseNet-121   & \texttt{densenet121.ra\_in1k}                & 256$\times$256 \\
InceptionV3    & \texttt{inception\_v3.tv\_in1k}              & 299$\times$299 \\
\bottomrule
\end{tabular}
\end{center}

The cyclone pre-training hyperparameters are swept across the following grid,
yielding 63 configurations in total (3 label-smoothing values $\times$ 3 learning
rates $\times$ 7 random seeds):

\begin{center}
\begin{tabular}{ll}
\toprule
\textbf{Parameter} & \textbf{Values} \\
\midrule
Cyclone label smoothing (\texttt{ls}) & 0.0, 0.1, 0.2 \\
Cyclone phase-2 learning rate (\texttt{lr}) & $10^{-6}$, $5\times10^{-6}$, $10^{-5}$ \\
Random seeds & 0, 1, 2, 3, 4, 5, 6 \\
\bottomrule
\end{tabular}
\end{center}

Bug-bite fine-tuning hyperparameters are fixed across all configurations:
label smoothing 0.0, phase-2 learning rate $5\times10^{-6}$.

% ============================================================
\subsection{Repository Structure}
\label{sec:pg-repo}
% ============================================================

\begin{description}
    \item[\texttt{notebooks/Multiclass\_Classification.ipynb}]
          The primary experimental notebook. Covers all four pipeline stages:
          control training, TC pre-training, feature map and GradCAM evaluation,
          and the full XAI suite. Intended for interactive exploration or
          one-shot reproduction on a GPU instance.

    \item[\texttt{Stacked\_Model\_Colab.ipynb}]
          Self-contained inference and visualisation notebook designed for
          Google Colab free tier. Loads trained weights from Google Drive,
          runs ensemble inference on a single image, and produces side-by-side
          Control vs.\ TC GradCAM and LIME outputs.

    \item[\texttt{scripts/}]
          Standalone Python scripts, one per pipeline stage. Intended for
          execution on a headless server via \texttt{run\_experiments.py}.
          \begin{itemize}
              \item \texttt{shared.py} --- shared imports, dataset paths,
                    metrics helpers, and GradCAM utilities used by all four scripts.
              \item \texttt{01\_train\_control.py} --- trains the Control
                    (ImageNet $\rightarrow$ bug-bite) models.
              \item \texttt{02\_train\_tc.py} --- runs cyclone pre-training then
                    TC $\rightarrow$ bug-bite fine-tuning.
              \item \texttt{03\_evaluate.py} --- generates feature map comparisons
                    and GradCAM overlays.
              \item \texttt{04\_xai.py} --- computes ensemble metrics, silhouette
                    scores, t-SNE projections, LIME, SHAP, and DiCE outputs.
          \end{itemize}

    \item[\texttt{miscellaneous\_code/pytorch\_utils.py}]
          Custom PyTorch training library imported by both the notebook and the
          scripts. See Section~\ref{sec:pg-pytorch-utils} for full documentation.

    \item[\texttt{miscellaneous\_code/run\_experiments.py}]
          Automated hyperparameter sweep runner. See
          Section~\ref{sec:pg-sweep} for full documentation.

    \item[\texttt{miscellaneous\_code/fetch\_cyclone\_dataset.py}]
          Downloads and labels the HURSAT-B1 cyclone dataset from NOAA.
          See Section~\ref{sec:pg-data-cyclone}.

    \item[\texttt{miscellaneous\_code/cyclone\_preprocessing.py}]
          Removes grid-line artefacts introduced by the HURSAT-B1 satellite
          image tiling process before training.

    \item[\texttt{Bug-Data/Multiclass\_Bug\_data/}]
          Bug-bite dataset, structured as \texttt{train/} and \texttt{val/}
          subdirectories, each containing one subdirectory per class.

    \item[\texttt{results/}]
          Auto-generated by the sweep runner. Contains checkpoints, per-run
          metric JSONs, XAI outputs, and plots.

    \item[\texttt{deploy/}]
          Legacy Streamlit demo (binary classifier, outdated). Not part of
          the active pipeline; retained for reference only.
\end{description}

% ============================================================
\subsection{Core Training Library: \texttt{pytorch\_utils.py}}
\label{sec:pg-pytorch-utils}
% ============================================================

This file contains all reusable PyTorch infrastructure. It is imported by
\texttt{run\_experiments.py}, by the four pipeline scripts in \texttt{scripts/},
and by the main notebook. Engineers extending the system should understand each
public function before modifying the training loop.

\subsubsection{\texttt{get\_pytorch\_loaders}}

\begin{lstlisting}[language=Python]
get_pytorch_loaders(train_dir, val_dir,
                    img_size=310, batch_size=4, augment=False)
\end{lstlisting}

Returns \texttt{(train\_loader, val\_loader, class\_names)} using
\texttt{torchvision.datasets.ImageFolder}. The \texttt{augment} parameter
accepts three values:

\begin{description}
    \item[\texttt{False}] No augmentation. Used for validation and for
          bug-bite fine-tuning in the sweep (the bug-bite dataset is small;
          augmentation during fine-tuning degraded performance in preliminary
          experiments).
    \item[\texttt{True}] Light augmentation: horizontal flip, $\pm10^{\circ}$
          rotation. Used for moderate augmentation scenarios.
    \item[\texttt{'strong'}] Aggressive augmentation designed for cyclone
          satellite imagery: random resized crop (scale 0.35--1.0), full
          $360^{\circ}$ rotation, vertical flip, colour jitter, random greyscale,
          Gaussian blur (kernel 5, $\sigma \in [0.5, 3.0]$), and random erasing.
          The full rotation and vertical flip are valid for cyclone images because
          satellite imagery is radially symmetric with respect to the storm centre.
\end{description}

All loaders use ImageNet normalisation
(\texttt{mean=[0.485, 0.456, 0.406]}, \texttt{std=[0.229, 0.224, 0.225]}).

\subsubsection{\texttt{train\_pytorch\_model}}

\begin{lstlisting}[language=Python]
train_pytorch_model(model, train_loader, val_loader, device,
                    phase1_epochs=5, phase2_epochs=20, patience=5,
                    save_path=None,
                    phase1_batch_size=None, phase2_batch_size=None,
                    label_smoothing=0.0, auto_class_weights=False,
                    phase1_lr=1e-4, phase2_lr=5e-6)
\end{lstlisting}

Implements the standard two-phase fine-tuning strategy used throughout the
project:

\begin{description}
    \item[Phase 1 (head-only)] All backbone parameters are frozen.
          Only the classification head (detected via \texttt{model.get\_classifier()})
          is trained for \texttt{phase1\_epochs} epochs at \texttt{phase1\_lr}
          with AdamW ($\lambda = 0.01$).
    \item[Phase 2 (full fine-tune)] All parameters are unfrozen and the
          entire network is trained at \texttt{phase2\_lr}. Gradient checkpointing
          is enabled where the model supports it.
\end{description}

Both phases use:
\begin{itemize}
    \item Early stopping on validation loss with patience \texttt{patience}.
    \item Automatic Mixed Precision (AMP) with \texttt{bfloat16} when a CUDA
          device is available.
    \item \texttt{CrossEntropyLoss} with optional label smoothing.
    \item Patience counters that reset independently between phases,
          because phase 2 typically causes an initial loss spike after
          unfreezing the backbone.
\end{itemize}

The best checkpoint (lowest validation loss) is saved to \texttt{save\_path}
and automatically reloaded into the model before the function returns.

\subsubsection{\texttt{get\_backbone\_state\_dict}}

\begin{lstlisting}[language=Python]
get_backbone_state_dict(state_dict)
\end{lstlisting}

Strips all classifier head keys from a state dict, enabling backbone weight
transfer between tasks with different numbers of output classes. Keys whose
names start with \texttt{head.}, \texttt{fc.}, or \texttt{classifier.} are
removed. This function is the mechanism by which the cyclone-pre-trained
backbone is transplanted into a fresh bug-bite classification head.

\subsubsection{Supporting Functions}

\begin{description}
    \item[\texttt{get\_pt\_probs\_preds(model, loader, device)}]
          Returns \texttt{(true\_labels, argmax\_preds, softmax\_probs)} as
          Python lists. Used for evaluation after training.
    \item[\texttt{evaluate\_pytorch\_model(model, val\_loader, class\_names, device)}]
          Prints a full \texttt{sklearn} classification report and renders a
          confusion matrix.
    \item[\texttt{plot\_pytorch\_history(history)}]
          Renders train/validation accuracy and loss curves from the history
          dict returned by \texttt{train\_pytorch\_model}.
\end{description}

% ============================================================
\subsection{Dataset Acquisition}
\label{sec:pg-data}
% ============================================================

\subsubsection{Bug-Bite Dataset}
\label{sec:pg-data-bug}

The bug-bite dataset contains images of five species categories:
\texttt{ants}, \texttt{bed\_bugs}, \texttt{mosquitos}, \texttt{spiders},
and \texttt{ticks\_fleas}.
It is organised as \texttt{train/\{class\}/} and \texttt{val/\{class\}/}
subdirectories inside \texttt{Bug-Data/Multiclass\_Bug\_data/}.

The dataset is available via two channels:
\begin{enumerate}
    \item \textbf{Repository:} the \texttt{Bug-Data/} directory is included in
          the repository and can be used directly.
    \item \textbf{Pre-processed archive:} a curated and pre-processed version,
          including augmented samples used in the published experiments, is
          available on request from\\\texttt{ishay.yulzary@gmail.com}.
\end{enumerate}

Dataset paths are configured at the top of \texttt{miscellaneous\_code/run\_experiments.py}
(\texttt{BUG\_TRAIN}, \texttt{BUG\_VAL}) and in \texttt{scripts/shared.py}.
They default to WSL native paths (\texttt{/home/test/bug\_data/\{train,val\}}),
which must be updated for other environments.

\subsubsection{Cyclone Dataset}
\label{sec:pg-data-cyclone}

The cyclone dataset is derived from two public sources:
\begin{itemize}
    \item \textbf{HURSAT-B1} (Hurricane Satellite, Band 1): infrared brightness
          temperature satellite imagery at $8\,\text{km}$ resolution,
          hosted by NOAA.
    \item \textbf{IBTrACS} (International Best Track Archive for Climate
          Stewardship): per-observation wind speed records used to assign
          intensity labels to each satellite frame.
\end{itemize}

Five intensity categories are defined, mapped to the Saffir-Simpson scale:

\begin{center}
\begin{tabular}{lll}
\toprule
\textbf{Category label} & \textbf{Wind speed (kt)} & \textbf{Saffir-Simpson equivalent} \\
\midrule
\texttt{cat1\_tropical\_depression} & 0--33   & Tropical Depression \\
\texttt{cat2\_tropical\_storm}      & 34--63  & Tropical Storm \\
\texttt{cat3\_hurricane\_cat12}     & 64--95  & Hurricane Cat 1--2 \\
\texttt{cat4\_hurricane\_cat34}     & 96--136 & Hurricane Cat 3--4 \\
\texttt{cat5\_hurricane\_cat5}      & 137+    & Hurricane Cat 5 \\
\bottomrule
\end{tabular}
\end{center}

Images cover the years 1995--2016. Each frame is clipped to the brightness
temperature range $[180, 300]\,\text{K}$, inverted (so cold cloud tops appear
bright), and saved as a $224\times224$ PNG.

\subsubsection*{Fetching the dataset}

\begin{lstlisting}[language=bash]
# Full dataset: 4000 images per category (20000 total, ~4h on a stable connection)
python miscellaneous_code/fetch_cyclone_dataset.py

# Quick trial run: 5 images per category
python miscellaneous_code/fetch_cyclone_dataset.py --trial

# Custom target size and worker count
python miscellaneous_code/fetch_cyclone_dataset.py --target 2000 --workers 8
\end{lstlisting}

The script is resumable: interrupted runs pick up from where they stopped
using \\\texttt{cyclone\_dataset/download\_log.json}. The output directory
defaults to \texttt{cyclone\_dataset/} in the working directory.

After fetching, the dataset must be split into train/val and placed at the
paths expected by the sweep runner. The pre-built split (500 val images per
category) is available from \texttt{ishay.yulzary@gmail.com}.

% ============================================================
\subsection{Training Pipeline: Individual Scripts}
\label{sec:pg-scripts}
% ============================================================

The four scripts in \texttt{scripts/} mirror the four sections of the main
notebook. They accept the same arguments and produce the same outputs, but
are designed for unattended server execution.

\subsubsection{\texttt{01\_train\_control.py}}

Trains the three Control models (ConvNeXt-Tiny, DenseNet-121, InceptionV3)
for one seed. Saves checkpoint files at the paths supplied via CLI arguments.

\begin{lstlisting}[language=bash]
python scripts/01_train_control.py \
    --run-id s0_ls0.10_lr5.0e-06 \
    --seed 0 \
    --ctrl-convnext  results/checkpoints/control_seed_0/control_convnext.pt \
    --ctrl-densenet  results/checkpoints/control_seed_0/control_densenet.pt \
    --ctrl-inception results/checkpoints/control_seed_0/control_inception.pt
\end{lstlisting}

\subsubsection{\texttt{02\_train\_tc.py}}

For one configuration, runs cyclone pre-training on all three architectures
and then fine-tunes each cyclone backbone on the bug-bite task.
Requires both cyclone checkpoint paths (output) and TC checkpoint paths (output).

\begin{lstlisting}[language=bash]
python scripts/02_train_tc.py \
    --run-id s0_ls0.10_lr5.0e-06 \
    --seed 0 \
    --label-smoothing 0.1 \
    --phase2-lr 5e-6 \
    --cyc-convnext  results/checkpoints/config_000/cyclone_convnext.pt \
    --cyc-densenet  results/checkpoints/config_000/cyclone_densenet.pt \
    --cyc-inception results/checkpoints/config_000/cyclone_inception.pt \
    --tc-convnext   results/checkpoints/config_000/tc_bug_convnext.pt \
    --tc-densenet   results/checkpoints/config_000/tc_bug_densenet.pt \
    --tc-inception  results/checkpoints/config_000/tc_bug_inception.pt
\end{lstlisting}

\subsubsection{\texttt{03\_evaluate.py}}

Generates feature map visualisations and GradCAM overlays comparing the
Control backbone, the TC cyclone backbone (before bug-bite fine-tuning),
and the TC fine-tuned model. Outputs are saved to
\texttt{results/feature\_maps/\{run-id\}/} and
\texttt{results/plots/\{run-id\}/}.

\subsubsection{\texttt{04\_xai.py}}

Computes the full XAI suite for one configuration:
ensemble accuracy (Control vs.\ TC), silhouette scores on backbone embeddings,
t-SNE projections, LIME superpixel explanations, SHAP gradient attributions,
and DiCE counterfactual explanations. Outputs are saved to
\texttt{results/xai/\{run-id\}/}.
Optional XAI libraries (\texttt{shap}, \texttt{lime}, \texttt{dice-ml}) are
checked at runtime; cells that depend on missing libraries are skipped
gracefully.

% ============================================================
\subsection{Automated Sweep Runner: \texttt{run\_experiments.py}}
\label{sec:pg-sweep}
% ============================================================

\texttt{miscellaneous\_code/run\_experiments.py} is the primary entry point for
training. It iterates over the full 63-configuration hyperparameter grid, calls
the four pipeline scripts in sequence for each configuration, and manages
checkpointing, state, and results atomically.

\subsubsection{Basic Usage}

\begin{lstlisting}[language=bash]
# Run all pending configurations
python miscellaneous_code/run_experiments.py

# Check progress without running
python miscellaneous_code/run_experiments.py --status

# Wipe state and results (checkpoints are preserved)
python miscellaneous_code/run_experiments.py --reset
\end{lstlisting}

\subsubsection{Configuration Grid}

The grid is defined by three constants at the top of the file:

\begin{lstlisting}[language=Python]
LABEL_SMOOTHING_VALUES = [0.0, 0.1, 0.2]
PHASE2_LR_VALUES       = [1e-6, 5e-6, 1e-5]
N_SEEDS                = 7
\end{lstlisting}

Adding a value to \texttt{LABEL\_SMOOTHING\_VALUES} or \texttt{PHASE2\_LR\_VALUES}
and re-running automatically appends the new configurations to the existing
state file without invalidating completed work.

\subsubsection{Pause and Resume}

The sweep can be interrupted and resumed at any point:

\begin{description}
    \item[Graceful pause between models:] create the file
          \texttt{results/PAUSE} (\texttt{touch results/PAUSE}) or press
          Ctrl-C. The runner finishes the current model, saves its checkpoint,
          and exits cleanly.
    \item[Crash recovery:] any configuration marked \texttt{running} in the
          state file (indicating a previous interrupted run) is automatically
          reset to \texttt{pending}. Already-saved checkpoints are reused,
          so only the incomplete model needs to be retrained.
\end{description}

\subsubsection{Output Files}

\begin{description}
    \item[\texttt{results/experiment\_state.json}]
          Tracks the status of every configuration (\texttt{pending},
          \texttt{running}, or \texttt{completed}).
    \item[\texttt{results/experiment\_results.json}]
          Per-model, per-configuration evaluation results: validation accuracy,
          macro F1, precision, recall, and full classification report.
    \item[\texttt{results/checkpoints/config\_\{id\}/}]
          TC cyclone checkpoint (\texttt{cyclone\_\{model\}.pt}) and
          TC bug-bite checkpoint (\texttt{tc\_bug\_\{model\}.pt}) for each
          configuration.
    \item[\texttt{results/checkpoints/control\_seed\_\{s\}/}]
          Control checkpoints, keyed by seed only (the Control pipeline is
          independent of the cyclone hyperparameters).
    \item[\texttt{results/xai/\{run-id\}/}]
          Per-configuration XAI outputs: \texttt{ensemble\_metrics.json},
          \\\texttt{silhouette.json}, t-SNE plots, GradCAM overlays,
          LIME visualisations, SHAP attribution maps.
\end{description}

\subsubsection{Hardware Requirements}

The full 63-configuration sweep was run on an NVIDIA RTX A5000 (24\,GB VRAM).
A GPU with at least 16\,GB VRAM is recommended for training. Estimated runtime
per configuration is 1.5--2 hours on equivalent hardware.
Inference and visualisation (scripts \texttt{03} and \texttt{04}) can be run
on CPU or a smaller GPU.

% ============================================================
\subsection{Experimental Results}
\label{sec:pg-results}
% ============================================================

\subsubsection{Ensemble Accuracy: Control vs.\ TC}

The table below reports the mean ensemble validation accuracy across all seven
seeds for each hyperparameter configuration. The ensemble averages softmax
probabilities across all three architectures before taking the argmax.

\begin{center}
\begin{tabular}{lcrrrr}
\toprule
\textbf{Config (ls, lr)} & \textbf{N seeds} &
\textbf{Control mean} & \textbf{TC mean} & \textbf{$\Delta$} \\
\midrule
ls=0.10, lr=$5\times10^{-6}$ & 7 & 81.83\% & 83.27\% & +1.45 pp \\
ls=0.00, lr=$1\times10^{-6}$ & 7 & 81.83\% & 83.18\% & +1.36 pp \\
ls=0.20, lr=$1\times10^{-6}$ & 7 & 81.83\% & 83.18\% & +1.36 pp \\
ls=0.20, lr=$5\times10^{-6}$ & 7 & 81.83\% & 83.00\% & +1.18 pp \\
ls=0.20, lr=$1\times10^{-5}$ & 7 & 81.83\% & 82.82\% & +0.99 pp \\
ls=0.00, lr=$1\times10^{-5}$ & 7 & 81.83\% & 82.73\% & +0.90 pp \\
ls=0.00, lr=$5\times10^{-6}$ & 7 & 81.83\% & 82.73\% & +0.90 pp \\
ls=0.10, lr=$1\times10^{-6}$ & 7 & 81.83\% & 82.28\% & +0.45 pp \\
ls=0.10, lr=$1\times10^{-5}$ & 7 & 81.83\% & 81.83\% & +0.00 pp \\
\bottomrule
\end{tabular}
\end{center}

TC pre-training matches or exceeds the Control baseline in all nine
configurations. The best single-run result is \textbf{86.08\%} TC ensemble
accuracy (seed 4, ls=0.00, lr=$10^{-5}$), compared to 81.01\% Control for
the same seed and configuration.

\subsubsection{Backbone Embedding Quality: Silhouette Scores}

Silhouette scores are computed on backbone embeddings extracted from the
validation set (before the classification head). Higher scores indicate
better class separation in the feature space.

\begin{center}
\begin{tabular}{lrr}
\toprule
\textbf{Architecture} & \textbf{Control} & \textbf{TC} \\
\midrule
ConvNeXt-Tiny & 0.057 & 0.081 \\
DenseNet-121  & 0.027 & 0.030 \\
InceptionV3   & \text{-}0.003 & \text{-}0.017 \\
\bottomrule
\end{tabular}
\end{center}

(Values shown are from the representative seed 0, ls=0.10, lr=$5\times10^{-6}$
configuration. Per-configuration scores are available in
\texttt{results/xai/\{run-id\}/silhouette.json}.)

% ============================================================
\subsection{XAI Methods}
\label{sec:pg-xai}
% ============================================================

All XAI outputs are generated by \texttt{scripts/04\_xai.py} and by Section 4
of the main notebook.

\begin{description}
    \item[GradCAM] Class Activation Mapping via gradients. Implemented in
          \texttt{scripts/shared.py}. Target layers are architecture-specific:
          the final stage for ConvNeXt, \texttt{norm5} for DenseNet,
          and \texttt{Mixed\_7c} for InceptionV3. Produces side-by-side
          Control vs.\ TC attention overlays for direct visual comparison.

    \item[LIME] Superpixel-based local explanation using
          \texttt{lime.lime\_image.LimeImageExplainer}. The ensemble prediction
          function is used as the black-box model. Top-$N$ segments
          (default 12) are coloured by their contribution sign and weight,
          using a red--blue diverging palette. Computationally expensive at
          high \texttt{num\_samples} values; the notebook default is 1000
          samples per image.

    \item[SHAP] Pixel-level attribution via \texttt{shap.GradientExplainer}.
          A background distribution of 50 validation images is used.
          Attribution maps are averaged across the three ensemble models
          and overlaid on the input image.

    \item[DiCE] Counterfactual explanations in embedding space. Backbone
          embeddings are reduced with PCA (50 components), and a surrogate
          logistic regression classifier is fitted on the validation set.
          DiCE generates counterfactual embedding vectors showing the minimal
          directional change required to flip the predicted class.

    \item[t-SNE] 2D projection of backbone embeddings for all five bug-bite
          classes, displayed as a $3 \times 2$ grid (one row per architecture,
          columns for Control and TC). Used to visualise whether TC pre-training
          produces more separable cluster structure.

    \item[Ensemble comparison] Soft-vote ensemble across all three architectures,
          evaluated on the validation set. Control and TC ensembles are compared
          directly; results saved to \texttt{ensemble\_metrics.json}.
\end{description}

% ============================================================
\subsection{Demo Notebook: \texttt{Stacked\_Model\_Colab.ipynb}}
\label{sec:pg-demo}
% ============================================================

This notebook provides a self-contained inference and visualisation interface
designed to run on Google Colab free tier. It does not require a local GPU or
local dataset.

\paragraph{Prerequisites.}
\begin{enumerate}
    \item A Google account with access to the project Google Drive folder.
          The folder must contain \texttt{project\_book\_results.tar},
          which packages all trained weight files.
    \item The notebook set \texttt{PROJECT\_ROOT} to the correct Drive path
          \\(default: \texttt{/content/drive/MyDrive/Capstone}).
\end{enumerate}

\paragraph{Workflow.}
\begin{enumerate}
    \item Mount Google Drive (Cell 2).
    \item Extract weights from the archive if not already extracted (Cell 3).
    \item Load all six models: three Control and three TC (Cell 5).
    \item Set \texttt{IMG\_PATH} to any JPEG or PNG image (Cell 7).
    \item Run inference: Cell 7 displays prediction probabilities for both
          pipelines side-by-side.
    \item Run GradCAM: Cell 8 renders a $3 \times 2$ overlay grid
          (architectures $\times$ pipelines).
    \item Run LIME: Cell 9 renders superpixel attributions for both pipelines.
\end{enumerate}

Weight file locations within the extracted archive:
\begin{itemize}
    \item Control: \texttt{project\_book\_results/s3\_ls0.20\_lr1.0e-05/}
    \item TC: \texttt{project\_book\_results/s6\_ls0.10\_lr5.0e-06/}
\end{itemize}

% ============================================================
\subsection{Extending the System}
\label{sec:pg-extend}
% ============================================================

\subsubsection{Adding a New Architecture}

\begin{enumerate}
    \item Add an entry to the \texttt{MODELS} list in
          \texttt{miscellaneous\_code/run\_experiments.py}:
          \texttt{(key, timm\_id, img\_size, cyc\_p1\_bs, bug\_bs,
          cyc\_p2\_epochs, bug\_p2\_epochs)}.
    \item Identify the correct GradCAM target layer for the new architecture
          and add a case to the \texttt{\_get\_target} function in
          \texttt{Stacked\_Model\_Colab.ipynb} and in \texttt{scripts/shared.py}.
    \item Re-run the sweep. New configurations are appended automatically;
          existing results are not invalidated.
\end{enumerate}

\subsubsection{Adding a New Bug-Bite Class}

\begin{enumerate}
    \item Add labelled images to \texttt{Bug-Data/Multiclass\_Bug\_data/train/}
          and \texttt{val/} under a new subdirectory named after the class.
    \item Update \texttt{N\_CLASSES} in \texttt{run\_experiments.py} if
          necessary (currently 5; \texttt{timm} models are instantiated
          with this value).
    \item Retrain from scratch. The existing Control and TC checkpoints are
          incompatible with a different number of output classes.
\end{enumerate}

\subsubsection{Adapting to a New Intermediate Domain}

The TC pipeline is not specific to cyclone imagery. Any image classification
task can serve as the intermediate pre-training domain by:
\begin{enumerate}
    \item Placing the new domain dataset at \texttt{TC\_TRAIN} / \texttt{TC\_VAL}
          in \texttt{run\_experiments.py}.
    \item Adjusting the augmentation strategy in \texttt{get\_pytorch\_loaders}
          if the new domain does not share cyclone imagery's symmetry properties.
    \item Updating the \texttt{augment='strong'} branch of
          \texttt{pytorch\_utils.py} accordingly.
\end{enumerate}

\subsubsection{Path Toward Productisation}

The current system is a research prototype. The following steps are required
before deployment:

\begin{enumerate}
    \item \textbf{Model serving:} wrap the ensemble inference logic from
          \texttt{Stacked\_Model\_Colab.ipynb} in a REST API endpoint
          (e.g.\ FastAPI). The core inference function is approximately
          30 lines of PyTorch code.
    \item \textbf{Weight hosting:} move trained weights from Google Drive to
          a persistent object store (e.g.\ AWS S3, GCP Cloud Storage).
    \item \textbf{Input validation:} add image size checks, colour space
          normalisation, and rejection of non-photographic inputs.
    \item \textbf{Dataset expansion:} the current bug-bite dataset is small
          (fewer than 500 images per class). Classification performance will
          improve substantially with a larger, clinically curated dataset.
    \item \textbf{Clinical validation:} model outputs should be reviewed and
          validated by a medical professional before any medical application.
\end{enumerate}

% --- End of appendix (remove if standalone) ---
\end{document}
